\documentclass[11pt]{article}
\usepackage[margin=1in]{geometry}

% Core packages
\usepackage{amsmath,amssymb}
\usepackage{multicol}
\usepackage{array}

% Paragraphs
\setlength{\parindent}{0pt}
\setlength{\parskip}{1\baselineskip}

\title{ChIP normalization}
\author{Benjamin Yeh}
\date{\today}

\begin{document}
\maketitle

\section{Introduction}

\section{Simple model}

Consider cells from two \textbf{conditons}: one control (vehicle), and one treated with a drug that decreases the abundance of a target epitope. Assume equal amounts of cells, antibody, and spike-in (exogenous) chromatin are used for IP across conditions. Consider the molar ratio in the input of the endogenous target epitope ($A$), the exogenous target epitope from the spike-in ($B$), off-target endogenous epitopes ($C$), and off-target exogenous epitopes from the spike-in ($D$):
\begin{itemize}
\item Control conditon: $A$:$B$:$C$:$D$
\item Treatment condition: $A'$:$B$:$C$:$D$
\end{itemize}

Assume that antibody-epitope binding reaches equilibrium with dissociation constants of $K_{d1}$ and $K_{d2}$ for target and off-target antibody-epitope interactions, respectively, so $K_{d1} \ll K_{d2}$.

Let $Y$ be antibody concentration. Let $Y_f$ refer to free antibody, while $Y_A$, $Y_B$, $Y_C$, and $Y_D$ refer to antibody-antigen complexes.

Key question: \textbf{What are the molar ratios of the epitopes after immunoprecipitation under different regimes of antibody concentration? What is the ratio of spike-in to endogenous epitopes in the immunoprecipitation?}

\subsection{Basic equations}

Because the endogenous and spike-in target epitopes have the same dissociation constant, the mass-action relation applies to each target population separately:

$$
\begin{gathered}
K_{d1}
= \frac{Y_f A_f}{Y_A}
= \frac{Y_f B_f}{Y_B}
= \frac{(Y_f) (A_f + B_f)}{Y_A + Y_B} \\
K_{d2}
= \frac{Y_f C_f}{Y_C}
= \frac{Y_f D_f}{Y_D}
= \frac{(Y_f) (C_f + D_f)}{Y_C + Y_D}
\end{gathered}
$$

Below, we will take advantage of the ratios

$$
\frac{Y_A}{A_f}=\frac{Y_f}{K_{d1}},
\qquad
\frac{Y_B}{B_f}=\frac{Y_f}{K_{d1}}.
$$

Conservation of mass:

$$\begin{gathered}
A = A_f + Y_A \\
B = B_f + Y_B \\
C = C_f + Y_C \\
D = D_f + Y_D \\
Y = Y_f + Y_A + Y_B + Y_C + Y_D
\end{gathered}$$

These equations can be used to derive expressions for the bound fractions of each class in terms of the concentration of free antibody $Y_f$ and the dissociation constants $K_{d1}$ and $K_{d2}$.

$$\begin{aligned}
\frac{Y_A}{A}
&= \frac{Y_A}{A_f + Y_A} \\
&= \frac{Y_A / A_f}{1 + Y_A / A_f} \\
&= \frac{Y_f / K_{d1}}{1 + Y_f / K_{d1}} \\
&= \frac{Y_f}{Y_f+K_{d1}}
\end{aligned}$$

By the same logic,
$$\begin{gathered}
\frac{Y_B}{B} = \frac{Y_f}{Y_f + K_{d1}} \\
\frac{Y_C}{C} = \frac{Y_f}{Y_f + K_{d2}} \\
\frac{Y_D}{D} = \frac{Y_f}{Y_f + K_{d2}} \\
\end{gathered}$$

In summary, the proportion bound of each epitope class is

$$\begin{gathered}
\theta_1
= \frac{Y_A}{A}
= \frac{Y_B}{B}
= \frac{Y_f}{Y_f+K_{d1}} \\
\theta_2
= \frac{Y_C}{C}
= \frac{Y_D}{D}
= \frac{Y_f}{Y_f+K_{d2}}
\end{gathered}$$

$K_{d1}$ is therefore the concentration of free antibody at equilibrium when half of the target epitopes are bound. Similarly, $K_{d2}$ is the concentration of free antibody at equilibrium when half of the off-target epitopes are bound.

The immunoprecipitation contains amounts $A \theta_1$, $B \theta_1$, $C \theta_2$, and $D \theta_2$. We can simplify this ratio to $A$ : $B$ : $\rho C$ : $\rho D$, where

$$
\rho = \frac{\theta_2}{\theta_1} = \frac{Y_f + K_{d1}}{Y_f + K_{d2}}.
$$

$\rho$ is the ratio of off-target capture efficiency to target capture efficiency and ranges from $r = \frac{K_{d1}}{K_{d2}} \ll 1$ when antibody is limiting ($Y_f \ll K_{d1}$) to 1 when antibody is in excess ($Y_f \gg K_{d2}$). Below, we use $\rho$ to denote this ratio for the control condition and $\rho'$ to denote this ratio for the treatment condition.

Additional observations:
\begin{itemize}
\item The ratios $A$ : $B$ and $C$ : $D$ are always conserved from input to immunoprecipitation.
\item The ratio of endogenous epitope to spike-in epitope is 
$$
\frac{A + \rho C}{B + \rho D}.
$$
\item The fraction of the immunoprecipitation from the spike-in is
$$
\frac{B + \rho D}{A + B + \rho C + \rho D}.
$$
As $\rho$ increases from $r$ to 1, this spike-in fraction can either
increase or decrease, depending on the sign of $AD-BC$. To see this,
differentiate with respect to $\rho$ using the quotient rule:
$$\begin{aligned}
\frac{d}{d\rho}\left(\frac{B+\rho D}{A+B+\rho(C+D)}\right)
&=\frac{D\bigl(A+B+\rho(C+D)\bigr)-(B+\rho D)(C+D)}
{\bigl(A+B+\rho(C+D)\bigr)^2} \\
&=\frac{AD+BD+\rho CD+\rho D^2-BC-BD-\rho CD-\rho D^2}
{\bigl(A+B+\rho(C+D)\bigr)^2} \\
&=\frac{AD-BC}{\bigl(A+B+\rho(C+D)\bigr)^2}.
\end{aligned}$$
The denominator is positive and the numerator does not depend on $\rho$,
so the spike-in fraction changes monotonically in $\rho$:
\begin{itemize}
    \item If $AD>BC$, i.e., $D/B>C/A$, it increases.
    \item If $AD<BC$, i.e., $D/B<C/A$, it decreases.
    \item If $AD=BC$, it is constant.
\end{itemize}
In words, the spike-in fraction increases with $\rho$ if and only if the
spike-in has more off-target epitope per target epitope than the cells
do.

\end{itemize}

\textbf{Regimes.} Because $\rho=(Y_f+K_{d1})/(Y_f+K_{d2})$ depends only on
the free-antibody concentration $Y_f$, the regimes below are classified by
where the antibody goes. Antibody conservation,
$Y=Y_f+S\theta_1+(C+D)\theta_2$ with $S=A+B$ (or $A'+B$), splits the
antibody among three ``sinks'': solution, target epitopes (capacity $S$),
and off-target epitopes (capacity $C+D$). Each regime corresponds to a
different sink dominating or saturating:
\begin{itemize}
    \item Regime 1 (limiting antibody): target epitopes absorb nearly all
    antibody.
    \item Regime 2 (dilute epitopes): epitopes are too dilute to absorb
    much antibody, so $Y_f\approx Y$.
    \item Regime 3 (antibody in excess of target): target epitopes are
    saturated, and the excess antibody is free or bound to unsaturated
    off-target epitopes.
    \item Regime 4 (saturating antibody): target and off-target epitopes
    are both saturated.
\end{itemize}
The boundaries between these regimes (antibody--target equivalence,
off-target half-saturation, and $K_{d1}\sim S$) are treated in the
crossovers section. A general result proved there is that
$\rho'\ge\rho$ whenever $A'<A$, in every regime: less target leaves more
free antibody, which raises off-target capture relative to target capture.
The table summarizes the results; the last column is the ChIP-Rx bias
factor $(B+\rho D)/(B+\rho'D)$ derived in Regime 1, by which total
spike-in normalization multiplies the true fold change $a_i'/a_i$.

{\renewcommand{\arraystretch}{1.6}
\footnotesize
\begin{center}
\begin{tabular}{|>{\raggedright\arraybackslash}p{0.13\linewidth}|>{\raggedright\arraybackslash}p{0.19\linewidth}|>{\raggedright\arraybackslash}p{0.15\linewidth}|>{\raggedright\arraybackslash}p{0.11\linewidth}|>{\raggedright\arraybackslash}p{0.12\linewidth}|>{\raggedright\arraybackslash}p{0.13\linewidth}|}
\hline
\textbf{Regime} & \textbf{Conditions} (for $S=A+B$ and $A'+B$) & $\boldsymbol{\rho}$ \textbf{(control)} & $\boldsymbol{\rho'}$ \textbf{(treatment)} & $\boldsymbol{\rho'}$ \textbf{vs.} $\boldsymbol{\rho}$ \textbf{when} $\boldsymbol{A'<A}$ & \textbf{ChIP-Rx bias} \\
\hline
1. Limiting antibody
& $Y<S$; $K_{d1}Y\ll(S-Y)^2$; $rY\ll S-Y$; $r(C+D)Y\ll(S-Y)^2$
& $\approx r\dfrac{A+B}{A+B-Y}$
& $\approx r\dfrac{A'+B}{A'+B-Y}$
& $\rho'>\rho$; both $\to r$ if $Y\ll S$
& $<1$; $\approx1$ if $\rho'D\ll B$ \\
\hline
2. Dilute epitopes
& $S\ll K_{d1}$; $C+D\ll K_{d2}$; any $Y$
& $\approx\dfrac{Y+K_{d1}}{Y+K_{d2}}$
& $\approx\rho$
& $\rho'\approx\rho$ (relative difference $\le S_c/K_{d1}+(C+D)/K_{d2}$)
& $\approx1$, even if $\rho D\gtrsim B$ \\
\hline
3. Antibody in excess of target
& $Y>S$; $rS(K_{d2}+C+D)\ll(Y-S)^2$; $Y-S\ll K_{d2}+C+D$
& $\approx r+\dfrac{Y-A-B}{K_{d2}+C+D}$
& $\approx r+\dfrac{Y-A'-B}{K_{d2}+C+D}$
& $\rho'>\rho$; $\rho'-\rho\approx\dfrac{A-A'}{K_{d2}+C+D}$, not suppressed by $r$
& $<1$, possibly far from 1 (overstates the decrease) \\
\hline
4. Saturating antibody
& $Y-S-(C+D)\gg K_{d2}$
& $\approx1$
& $\approx1$
& $\rho'>\rho$, but negligibly
& $\approx1$, but no enrichment over off-target (IP $\approx$ input) \\
\hline
Crossover: equivalence (1$\leftrightarrow$3)
& $|S-Y|\lesssim$ a few $\sqrt{K_{d1}S}$; $K_{d1}\ll S$
& $\approx\dfrac{Y_f+K_{d1}}{K_{d2}}$, $Y_f\approx\sqrt{K_{d1}S}$ at $Y=S$
& same with $S=A'+B$
& $\rho'\gg\rho$ possible
& $<1$; largest bias of any regime \\
\hline
Crossover: off-target half-saturation (3$\leftrightarrow$4)
& Target saturated; $Y_f\sim K_{d2}$
& $\approx\dfrac{Y_f}{Y_f+K_{d2}}$; $=\tfrac12$ at $Y=S+K_{d2}+\tfrac{C+D}{2}$
& same with $S=A'+B$
& $\rho'>\rho$, slightly
& $\le1$; $\approx1$ unless $A-A'\sim K_{d2}$ \\
\hline
\end{tabular}
\end{center}
}

\subsection{Regime 1: limiting antibody}

Let $S_c=A+B$ and $S_t=A'+B$ be the total target-epitope
concentrations in the control and treatment conditions, and let
$S\in\{S_c,S_t\}$ denote the total target-epitope concentration of
either condition. We assume the following
for both conditions (i.e., for each $S\in\{S_c,S_t\}$).

The assumptions are easiest to interpret through a heuristic estimate of
the free-antibody concentration. If nearly all antibody is bound to
target, then about $Y$ target epitopes are bound and about $S-Y$ are free.
Rearranging the mass-action relation
$K_{d1}=Y_f(A_f+B_f)/(Y_A+Y_B)$ then gives
$$
Y_f
=K_{d1}\frac{Y_A+Y_B}{A_f+B_f}
\approx K_{d1}\frac{Y}{S-Y}.
$$
Below, we show that this estimate is in fact an exact upper bound on
$Y_f$.

\begin{enumerate}
    \item Antibody is less abundant than target epitope:
    $$
    Y<S.
    $$
    Thus, even if every antibody molecule binds a target epitope, the
    remaining target concentration $S-Y$ is positive.

    \item Free antibody is scarce relative to the remaining target:
    $$
    K_{d1}Y\ll(S-Y)^2.
    $$
    Dividing both sides by $S-Y$ shows that this is equivalent to
    $K_{d1}Y/(S-Y)\ll S-Y$: the estimated free antibody is much less than
    the free target. This matters because antibody that is free does not
    occupy a target epitope, so it leaves an extra epitope free; the free
    target concentration stays close to $S-Y$ only if $Y_f$ is small
    compared with $S-Y$. The square arises because $S-Y$ appears twice:
    once in the estimate of $Y_f$ (more free target pulls more antibody out
    of solution) and once as the yardstick against which $Y_f$ is compared.

    Taking square roots, the condition is equivalent to
    $\sqrt{K_{d1}Y}\ll S-Y$. Here $\sqrt{K_{d1}Y}$ is roughly the
    free-antibody concentration at exact equivalence ($S=Y$): there, the
    bound target is approximately $Y$ and the free target is
    $S-(Y-Y_f)=Y_f$, so the mass-action relation gives
    $K_{d1}\approx Y_f\cdot Y_f/Y$, i.e., $Y_f\approx\sqrt{K_{d1}Y}$. The
    condition thus says that the excess target $S-Y$ must greatly exceed
    the free antibody that would remain if the antibody exactly titrated the
    target.

    Formally, below we show that $Y_f\le K_{d1}Y/(S-Y)$. Dividing both
    sides by $S-Y$ (which is positive by assumption 1) and then applying
    this condition,
    $$
    \frac{Y_f}{S-Y}
    \le\frac{K_{d1}Y}{(S-Y)(S-Y)}
    =\frac{K_{d1}Y}{(S-Y)^2}
    \ll1.
    $$
    Multiplying both sides by $S-Y$ gives $Y_f\ll S-Y$.

    \item The free-antibody concentration is small relative to $K_{d2}$:
    $$
    rY\ll S-Y.
    $$
    Dividing both sides by $S-Y$ and substituting $r=K_{d1}/K_{d2}$,
    $$
    \frac{rY}{S-Y}
    =\frac{1}{K_{d2}}\cdot K_{d1}\frac{Y}{S-Y}
    \ll1,
    $$
    so this condition says that the estimated free antibody is much less
    than $K_{d2}$: off-target epitopes are far from half-occupied. In this
    unsaturated regime, $\theta_2=Y_f/(Y_f+K_{d2})\approx Y_f/K_{d2}$, so
    $rY/(S-Y)$ is approximately the fraction of off-target epitopes that are
    bound.

    The quantity $rY$ is not especially meaningful on its own; the
    rearrangement $Y\ll(S-Y)/r$ is more interpretable. Rearranging the
    mass-action relation, the target-bound antibody is
    $Y_A+Y_B=Y_f(A_f+B_f)/K_{d1}$. For $Y_f$ to reach $K_{d2}$ while about
    $S-Y$ target remains free, the target-bound antibody would have to be
    about $K_{d2}(S-Y)/K_{d1}=(S-Y)/r$. The condition says that the total
    antibody is far below this amount. Below, we show that this condition
    implies $Y_f\ll K_{d2}$.

    \item Off-target binding consumes little antibody:
    $$
    r(C+D)Y\ll(S-Y)^2.
    $$
    By the mass-action relations, a free target epitope is bound by
    antibody in the ratio $(Y_A+Y_B)/(A_f+B_f)=Y_f/K_{d1}$, while an
    off-target epitope is bound in the ratio
    $(Y_C+Y_D)/(C_f+D_f)=Y_f/K_{d2}$. In the unsaturated regime
    (assumption 3), nearly all off-target epitopes are free, so
    $C_f+D_f\approx C+D$ and
    $$
    Y_C+Y_D
    \approx\frac{Y_f}{K_{d2}}(C+D)
    =\frac{Y_f}{K_{d1}}\cdot r(C+D).
    $$
    That is, the off-target pool binds as much antibody as $r(C+D)$ free
    target epitopes would. Thus $r(C+D)$ is the off-target pool measured in
    ``free-target-epitope equivalents.''

    Dividing both sides of the condition by $S-Y$ and substituting
    $r=K_{d1}/K_{d2}$,
    $$
    (C+D)\cdot\frac{1}{K_{d2}}\cdot K_{d1}\frac{Y}{S-Y}\ll S-Y.
    $$
    The left-hand side is $C+D$ times the estimated off-target occupancy
    $Y_f/K_{d2}$, i.e., the estimated off-target-bound antibody. As in
    assumption 2, the condition says that antibody that is not on a target
    epitope (here, antibody bound off-target) is much less than the free
    target, and the square arises for the same reason. Below, we show that
    this condition makes the off-target-bound antibody amount negligible
    relative to the remaining target $S-Y$.
\end{enumerate}
In summary, assumptions 2 and 4 keep the free target concentration close
to $S-Y$ by limiting the two places antibody can go other than a target
epitope: solution and off-target epitopes. Assumption 3 keeps off-target
binding unsaturated.

\textbf{Simpler sufficient conditions.} Suppose $Y\le S/2$. Multiplying
both sides by 2 gives $2Y\le S$, and subtracting $Y$ from both sides gives
$Y\le S-Y$. Since $S-Y>0$ by assumption 1, dividing both sides by $S-Y$
gives
$$
\frac{Y}{S-Y}\le1.
$$
Now write the left-hand sides of assumptions 2 and 4, divided by
$(S-Y)^2$, as products:
$$
\frac{K_{d1}Y}{(S-Y)^2}
=\frac{K_{d1}}{S-Y}\cdot\frac{Y}{S-Y}
\le\frac{K_{d1}}{S-Y},
\qquad
\frac{r(C+D)Y}{(S-Y)^2}
=\frac{r(C+D)}{S-Y}\cdot\frac{Y}{S-Y}
\le\frac{r(C+D)}{S-Y},
$$
where each inequality uses $Y/(S-Y)\le1$. Therefore, if
$K_{d1}\ll S-Y$, then $K_{d1}Y/(S-Y)^2\ll1$, which is assumption 2; and
if $r(C+D)\ll S-Y$, then $r(C+D)Y/(S-Y)^2\ll1$, which is assumption 4.
When $Y$ approaches $S$, however, the factor $Y/(S-Y)$ grows without
bound, and the simpler conditions are no longer sufficient.

We now derive $\rho$ for either condition. (For the treatment condition,
$\theta_1$, $\theta_2$, $Y_f$, and $\rho$ should be read as their primed
counterparts.)

\textbf{Antibody conservation.} The target-bound antibody is
$Y_A+Y_B=A\theta_1+B\theta_1=S\theta_1$, and the off-target-bound antibody
is $Y_C+Y_D=C\theta_2+D\theta_2=(C+D)\theta_2$. Substituting these into
$Y=Y_f+Y_A+Y_B+Y_C+Y_D$ gives
$$
Y=Y_f+S\theta_1+(C+D)\theta_2.
$$

\textbf{Step 1: free target.} The free target concentration is
$S-S\theta_1=S(1-\theta_1)$. Solving antibody conservation for
$S\theta_1$ gives $S\theta_1=Y-Y_f-(C+D)\theta_2$, so
\begin{equation}
S(1-\theta_1)
=S-\bigl(Y-Y_f-(C+D)\theta_2\bigr)
=(S-Y)+Y_f+(C+D)\theta_2.
\label{eq:r1-1}
\end{equation}
Intuitively, the free target is the target left over if all antibody
were bound to target ($S-Y$), plus the target ``freed up'' by antibody
that is free or bound off-target.

\textbf{Step 2: an exact expression for $\rho$.} From
$\theta_1=Y_f/(Y_f+K_{d1})$,
$$
1-\theta_1
=\frac{Y_f+K_{d1}}{Y_f+K_{d1}}-\frac{Y_f}{Y_f+K_{d1}}
=\frac{K_{d1}}{Y_f+K_{d1}},
$$
so, multiplying the numerator and denominator by $S$,
\begin{equation}
Y_f+K_{d1}=\frac{K_{d1}}{1-\theta_1}=\frac{K_{d1}S}{S(1-\theta_1)}.
\label{eq:r1-2}
\end{equation}
Substituting \eqref{eq:r1-2} into the numerator of $\rho$, and then \eqref{eq:r1-1} into the
result,
$$\begin{aligned}
\rho
&=\frac{Y_f+K_{d1}}{Y_f+K_{d2}} \\
&=\frac{K_{d1}S}{S(1-\theta_1)\,(Y_f+K_{d2})} \\
&=\frac{K_{d1}S}{\bigl[(S-Y)+Y_f+(C+D)\theta_2\bigr](K_{d2}+Y_f)}.
\end{aligned}$$
Factoring $S-Y$ out of the first bracket and $K_{d2}$ out of the second,
and using $r=K_{d1}/K_{d2}$,
\begin{equation}
\rho
=r\frac{S}{S-Y}\cdot\frac{1}{(1+\varepsilon_1)(1+\varepsilon_2)},
\qquad
\varepsilon_1=\frac{Y_f+(C+D)\theta_2}{S-Y},
\quad
\varepsilon_2=\frac{Y_f}{K_{d2}}.
\label{eq:r1-3}
\end{equation}
Equation \eqref{eq:r1-3} is exact. It remains to show that $\varepsilon_1\ll1$ and
$\varepsilon_2\ll1$.

\textbf{Step 3: an upper bound on $Y_f$.} Because $Y_f\ge0$ and
$(C+D)\theta_2\ge0$, equation \eqref{eq:r1-1} gives $S(1-\theta_1)\ge S-Y$, which is
positive by assumption 1. A smaller denominator gives a larger fraction,
so substituting into \eqref{eq:r1-2},
$$
Y_f+K_{d1}\le\frac{K_{d1}S}{S-Y}.
$$
Subtracting $K_{d1}$ from both sides,
\begin{equation}
Y_f
\le\frac{K_{d1}S}{S-Y}-\frac{K_{d1}(S-Y)}{S-Y}
=\frac{K_{d1}Y}{S-Y}.
\label{eq:r1-4}
\end{equation}

\textbf{Step 4: $\varepsilon_2\ll1$.} Dividing \eqref{eq:r1-4} by $K_{d2}$,
\begin{equation}
\varepsilon_2
=\frac{Y_f}{K_{d2}}
\le\frac{K_{d1}}{K_{d2}}\cdot\frac{Y}{S-Y}
=\frac{rY}{S-Y}
\ll1,
\label{eq:r1-5}
\end{equation}
where the final inequality is assumption 3. Therefore, $Y_f\ll K_{d2}$.

\textbf{Step 5: $\varepsilon_1\ll1$.} Split $\varepsilon_1$ into its two
terms. For the free-antibody term, dividing \eqref{eq:r1-4} by $S-Y$,
$$
\frac{Y_f}{S-Y}
\le\frac{K_{d1}Y}{(S-Y)^2}
\ll1,
$$
by assumption 2. For the off-target term, dropping $Y_f$ from the
denominator of $\theta_2$ makes it larger, so by \eqref{eq:r1-5},
$$
\theta_2=\frac{Y_f}{Y_f+K_{d2}}\le\frac{Y_f}{K_{d2}}\le\frac{rY}{S-Y}.
$$
Multiplying by $(C+D)/(S-Y)$,
$$
\frac{(C+D)\theta_2}{S-Y}
\le\frac{r(C+D)Y}{(S-Y)^2}
\ll1,
$$
by assumption 4. Since both terms are $\ll1$, so is their sum
$\varepsilon_1$.

\textbf{Step 6: result.} Since $\varepsilon_1\ll1$ and $\varepsilon_2\ll1$,
$(1+\varepsilon_1)(1+\varepsilon_2)\approx1$, and \eqref{eq:r1-3} becomes
$$
\rho\approx r\frac{S}{S-Y}.
$$
(Because $\varepsilon_1,\varepsilon_2\ge0$, \eqref{eq:r1-3} also shows that
$r\,S/(S-Y)$ is an exact upper bound on $\rho$.) Substituting $S=A+B$ for
the control and $S=A'+B$ for the treatment,
$$
\rho\approx r\frac{A+B}{A+B-Y},
\qquad
\rho'\approx r\frac{A'+B}{A'+B-Y}.
$$

If, in addition, $Y\ll\min(A+B,A'+B)$, these reduce to $\rho\approx\rho'\approx r$.

Note that \textbf{$\rho'$ increases relative to $\rho$ when $A' < A$}. This can be seen from the expression
$$\rho \approx r \frac{S}{S-Y} = r \frac{S-Y+Y}{S-Y} = r \left(1 + \frac{Y}{S-Y} \right).$$
For $S>Y$ (assumption 1), decreasing $S$ decreases the positive denominator $S-Y$, so $\frac{Y}{S-Y}$ and hence $\rho$ increase.

\subsubsection{ChIP normalization}

In the control condition, the ratio of endogenous to spike-in epitopes in the immunoprecipitation is
$\frac{A + \rho C}{B + \rho D}$

In the treatment condition, the ratio of endogenous to spike-in epitopes in the immunoprecipitation is
$\frac{A' + \rho' C}{B + \rho' D}$.

\textbf{Goal.} Let $a_i$ and $a_i'$ be the concentrations of target
epitope at genomic locus $i$ in the control and treatment conditions,
so that $\sum_i a_i=A$ and $\sum_i a_i'=A'$. We want normalized signals
whose ratio between conditions at each locus is the true fold change
$a_i'/a_i$.

\textbf{Reads.} Assume that sequencing samples reads in proportion to the
composition of the immunoprecipitation, and that target loci do not
overlap off-target loci. Let $N$ be the total number of reads in the
control library. The immunoprecipitation contains $a_i\theta_1$ at locus
$i$ out of a total of $(A+B)\theta_1+(C+D)\theta_2$. Dividing the numerator
and denominator by $\theta_1$ and substituting $\rho=\theta_2/\theta_1$,
the expected number of reads at locus $i$ is
$$
n_i
=N\frac{a_i\theta_1}{(A+B)\theta_1+(C+D)\theta_2}
=N\frac{a_i}{A+B+\rho(C+D)}.
$$
By the same logic, the expected number of reads mapping to the spike-in
genome (from both target and off-target spike-in epitopes) is
$$
N_s
=N\frac{B+\rho D}{A+B+\rho(C+D)}.
$$
The treatment library has the analogous quantities $N'$, $n_i'$, and
$N_s'$, with $A$, $a_i$, and $\rho$ replaced by $A'$, $a_i'$, and $\rho'$.

\textbf{Scaling by total reads (RPM).} Multiplying each library by
$\alpha=1/N$ (or $10^6/N$, in reads per million) gives
$$
\frac{n_i}{N}=\frac{a_i}{A+B+\rho(C+D)},
\qquad
\frac{n_i'}{N'}=\frac{a_i'}{A'+B+\rho'(C+D)}.
$$
The ratio between conditions is
$$
\frac{n_i'/N'}{n_i/N}
=\frac{a_i'}{a_i}\cdot\frac{A+B+\rho(C+D)}{A'+B+\rho'(C+D)}.
$$
If the treatment decreases the target globally ($A'<A$) and the
spike-in and off-target terms are small, the second factor is
approximately $A/A'>1$, which cancels the global decrease. Thus, RPM
normalization hides global changes; this is the motivation for spike-in
normalization.

\textbf{Scaling by spike-in reads (ChIP-Rx).} Multiplying each library by
$\alpha=1/N_s$ for the control and $\alpha'=1/N_s'$ for the treatment (in
practice, $10^6/N_s$ and $10^6/N_s'$), the factors $N$ and
$A+B+\rho(C+D)$ cancel:
$$
\frac{n_i}{N_s}
=\frac{N\,a_i/(A+B+\rho(C+D))}{N\,(B+\rho D)/(A+B+\rho(C+D))}
=\frac{a_i}{B+\rho D},
\qquad
\frac{n_i'}{N_s'}=\frac{a_i'}{B+\rho' D}.
$$
The ratio between conditions is
$$
\frac{n_i'/N_s'}{n_i/N_s}
=\frac{a_i'}{a_i}\cdot\frac{B+\rho D}{B+\rho' D}.
$$
The second factor is a bias. It equals 1 if $\rho D\ll B$ and
$\rho' D\ll B$, i.e., if off-target binding to the spike-in is
negligible compared with target binding to the spike-in. In the
limiting-antibody regime, substituting the approximations for $\rho$ and
$\rho'$, this requires
$$
r\frac{A+B}{A+B-Y}D\ll B,
\qquad
r\frac{A'+B}{A'+B-Y}D\ll B.
$$
Otherwise, because $\rho'>\rho$ when $A'<A$, we have $B+\rho' D>B+\rho D$,
so the bias factor is less than 1: spike-in normalization
\emph{overstates} the decrease in the treatment condition.

\textbf{Scaling by spike-in target reads.} The bias arises because $N_s$
includes off-target spike-in reads, whose capture efficiency differs
between conditions. If instead we count only the reads $N_{s,B}$ at the
spike-in's target loci (e.g., reads in spike-in peaks), then by the same
logic as above,
$$
N_{s,B}
=N\frac{B}{A+B+\rho(C+D)},
\qquad
N_{s,B}'
=N'\frac{B}{A'+B+\rho'(C+D)}.
$$
Using scaling factors $\alpha=1/N_{s,B}$ and $\alpha'=1/N_{s,B}'$,
$$
\frac{n_i}{N_{s,B}}=\frac{a_i}{B},
\qquad
\frac{n_i'}{N_{s,B}'}=\frac{a_i'}{B},
\qquad
\frac{n_i'/N_{s,B}'}{n_i/N_{s,B}}=\frac{a_i'}{a_i},
$$
which recovers the true fold change regardless of $\rho$ and $\rho'$.
(In practice, reads in spike-in peaks also include some off-target
background, so this cancellation is only approximate.)

\subsubsection{Example: Pol II S2ph ChIP in K562 cells with an mESC spike-in}

We now ask what antibody and spike-in amounts would place a realistic
experiment in Regime 1. All parameter values below are order-of-magnitude
estimates; the point is the method, which can be rerun with better
numbers.

\textbf{Experimental setup.}
\begin{itemize}
    \item 5 million K562 cells per condition; control is vehicle, and the
    treatment (e.g., a CDK9 inhibitor) halves Ser2-phosphorylated Pol II,
    so $A'=A/2$.
    \item Spike-in: mouse embryonic stem cell (mESC) chromatin from
    $5\times10^5$ cells (10\% of the K562 cell number).
    \item Antibody: Cell Signaling Technology \#13499, Phospho-Rpb1 CTD
    (Ser2) (E1Z3G) rabbit monoclonal antibody (IgG, $\approx150$ kDa).
    The manufacturer recommends 10~$\mu$L (1:50) per IP of 10~$\mu$g
    chromatin ($\approx4\times10^6$ cells). The stock concentration is
    lot-specific and not listed on the product page; we assume
    $50~\mu\text{g/mL}$, so 10~$\mu$L is $\approx0.5~\mu$g.
    \item IP volume $V=500~\mu$L (typical for SimpleChIP-style protocols).
\end{itemize}

\textbf{Epitope abundances.} A HeLa cell contains roughly
$3\times10^5$ Pol II molecules, of which roughly $0.65$--$1.1\times10^5$
are hyperphosphorylated (Pol IIO) or actively transcribing (Kimura et
al., 1999; Mol Cell Biol 19:5383). K562 cells are somewhat larger than
HeLa cells with substantially larger nuclei, so we take
$10^5$ Ser2-phosphorylated Pol II per K562 cell. We are not aware of a
direct count in mESCs; given their smaller size, we take
$5\times10^4$ per mESC. We treat each Ser2-phosphorylated Pol II as a
single target epitope, even though the CTD contains 52 heptad repeats
that can each carry Ser2 phosphorylation. Multiple epitopes per
molecule mainly increase the effective (avidity-enhanced) affinity, which
we absorb into $K_{d1}$. For off-target epitopes (non-S2ph Pol II, other
phosphoproteins, and nonspecific chromatin sites), we assume $10^7$ sites
per cell in both species.

To convert a number of molecules per cell into a concentration in the IP,
multiply by the number of cells, divide by Avogadro's number
$N_A=6.022\times10^{23}~\text{mol}^{-1}$, and divide by $V$:
$$\begin{aligned}
A
&=\frac{(5\times10^6)(10^5)}{(6.022\times10^{23})(5\times10^{-4}~\text{L})}
=\frac{5\times10^{11}}{3.011\times10^{20}}~\text{M}
\approx1.66~\text{nM}, \\
B
&=\frac{(5\times10^5)(5\times10^4)}{(6.022\times10^{23})(5\times10^{-4}~\text{L})}
\approx0.083~\text{nM}, \\
C
&=\frac{(5\times10^6)(10^7)}{(6.022\times10^{23})(5\times10^{-4}~\text{L})}
\approx166~\text{nM},
\qquad
D\approx16.6~\text{nM}.
\end{aligned}$$
With $A'=A/2\approx0.83$~nM, the total target concentrations are
$$
S_c=A+B\approx1.74~\text{nM},
\qquad
S_t=A'+B\approx0.91~\text{nM},
\qquad
C+D\approx183~\text{nM}.
$$

\textbf{Affinities.} We assume an effective target affinity
$K_{d1}=0.1$~nM (typical of rabbit monoclonal antibodies, and helped by
avidity) and an off-target affinity $K_{d2}=1~\mu\text{M}=1000$~nM, so
$r=K_{d1}/K_{d2}=10^{-4}$.

\textbf{The recommended antibody amount is not in Regime 1.} An IgG of
150~kDa has $1/(1.5\times10^5~\text{g/mol})\approx6.7$~pmol per $\mu$g.
Thus $0.5~\mu$g in 500~$\mu$L is
$$
Y
=\frac{0.5\times10^{-6}~\text{g}}{(1.5\times10^5~\text{g/mol})(5\times10^{-4}~\text{L})}
\approx6.7~\text{nM},
$$
which exceeds both $S_c$ and $S_t$, violating assumption 1. The
recommended protocol therefore operates with antibody in excess of
target (Regime 3 below).

\textbf{Choosing $Y$ for Regime 1.} The treatment condition, with the
smaller $S$, is the binding constraint, so we check each assumption with
$S=S_t\approx0.91$~nM. Try $Y=0.2$~nM, so $S_t-Y\approx0.71$~nM.
\begin{enumerate}
    \item $Y=0.2~\text{nM}<S_t\approx0.91$~nM. \checkmark
    \item $\dfrac{K_{d1}Y}{(S_t-Y)^2}
    =\dfrac{(0.1)(0.2)}{(0.71)^2}
    \approx0.04\ll1$. \checkmark
    \item $\dfrac{rY}{S_t-Y}
    =\dfrac{(10^{-4})(0.2)}{0.71}
    \approx3\times10^{-5}\ll1$. \checkmark
    \item $\dfrac{r(C+D)Y}{(S_t-Y)^2}
    =\dfrac{(10^{-4})(183)(0.2)}{(0.71)^2}
    \approx0.007\ll1$. \checkmark
\end{enumerate}
In the control, $S_c-Y\approx1.54$~nM is larger, so every ratio is even
smaller. Assumption 2 is the tightest: increasing $Y$ to 0.45~nM
($\approx S_t/2$) would raise its ratio to
$(0.1)(0.45)/(0.46)^2\approx0.2$, which is no longer $\ll1$.

Converting $Y=0.2$~nM into an amount of antibody,
$$
(0.2\times10^{-9}~\text{M})(5\times10^{-4}~\text{L})
=10^{-13}~\text{mol},
\qquad
(10^{-13}~\text{mol})(1.5\times10^5~\text{g/mol})
=1.5\times10^{-8}~\text{g}
=15~\text{ng},
$$
i.e., about 0.3~$\mu$L of the assumed 50~$\mu$g/mL stock, roughly
30-fold less than recommended. (If both arms of the IgG bind
independently, the binding-site concentration is $2Y$, and the antibody
amount should be halved again.)

\textbf{Predicted capture efficiencies and bias.} Substituting into the
Regime 1 result,
$$\begin{aligned}
\rho&\approx r\frac{S_c}{S_c-Y}=10^{-4}\cdot\frac{1.74}{1.54}\approx1.13\times10^{-4}, \\
\rho'&\approx r\frac{S_t}{S_t-Y}=10^{-4}\cdot\frac{0.91}{0.71}\approx1.28\times10^{-4}.
\end{aligned}$$
The ChIP-Rx bias factor is then
$$
\frac{B+\rho D}{B+\rho'D}
=\frac{0.083+(1.13\times10^{-4})(16.6)}{0.083+(1.28\times10^{-4})(16.6)}
=\frac{0.0849}{0.0851}
\approx0.997,
$$
so total-spike-in normalization is accurate to $\approx0.3\%$. For
comparison, solving the full conservation equation numerically at the
recommended $Y\approx6.7$~nM gives $\rho\approx4.3\times10^{-3}$,
$\rho'\approx5.0\times10^{-3}$, and a bias factor of
$\approx0.93$: ChIP-Rx would overstate the treatment-induced decrease by
$\approx7\%$, and off-target epitopes would make up nearly half
($\rho D/(B+\rho D)\approx0.46$) of the spike-in reads.

\textbf{Yield.} In Regime 1 nearly all antibody is bound to target, so
about $Y\cdot V=10^{-13}$~mol $\approx6\times10^{10}$ Pol II molecules are
captured: $\theta_1\approx Y/S$ is about 11\% of target epitopes in the
control and 22\% in the treatment. At $\sim$200~bp per captured fragment
and 650~g/mol per base pair, this is roughly
$(10^{-13})(200)(650)\approx13$~ng of DNA, before losses, which is ample
for library preparation.

\textbf{Practical notes.} If the antibody amount cannot be reduced this
far, increasing $S$ helps: concentrating the IP into a smaller volume or
using more cells raises $S-Y$ for a fixed amount of antibody, which
relaxes assumption 2 quadratically. Because the spike-in contributes
only $\approx5\%$ of the target ($B/A\approx0.05$), changing the spike-in
amount barely affects whether Regime 1 holds. Its main role is to
ensure enough spike-in target reads, and to keep $\rho D\ll B$. Finally,
the stock concentration, $K_{d1}$, and the epitope counts should be
measured or checked (e.g., against the lot's certificate of analysis) before
relying on these numbers.

\subsection{Regime 2: dilute epitopes (no antibody depletion)}

As in Regime 1, let $S_c=A+B$ and $S_t=A'+B$ be the total target-epitope
concentrations in the control and treatment conditions, and let
$S\in\{S_c,S_t\}$ denote either one. In Regime 1, antibody was the scarce
reagent: nearly every antibody molecule found a target, so the amount of
target captured was fixed at about $Y$ and the \emph{fraction} of target
captured depended on $S$. Here we consider the opposite situation: the
epitopes are so dilute that binding them removes only a negligible amount
of antibody from solution. The free antibody concentration is then set
by the experimenter ($Y_f\approx Y$) rather than by the chromatin, and,
as we will see, every capture efficiency becomes independent of $S$.

We assume the following for both conditions (i.e., for each
$S\in\{S_c,S_t\}$).

\begin{enumerate}
    \item Target epitopes are dilute relative to $K_{d1}$:
    $$
    S\ll K_{d1}.
    $$
    To interpret the ratio $S/K_{d1}$, rearrange the mass-action relation
    $K_{d1}=Y_f(A_f+B_f)/(Y_A+Y_B)$ to give the ratio of target-bound to
    free antibody:
    $$
    \frac{Y_A+Y_B}{Y_f}=\frac{A_f+B_f}{K_{d1}}.
    $$
    When the target is unsaturated, nearly all target is free
    ($A_f+B_f\approx S$), so $S/K_{d1}$ is approximately the ratio of
    target-bound antibody to free antibody. The assumption therefore says
    that for every antibody molecule sitting on a target epitope, many more
    are free in solution. (If the target is partly saturated, fewer target
    epitopes are free, $A_f+B_f<S$, so $S/K_{d1}$ is an upper bound on this
    ratio.) Note the contrast with Regime 1, which required the
    opposite inequality, $K_{d1}\ll S-Y$.

    \item Off-target epitopes are dilute relative to $K_{d2}$:
    $$
    C+D\ll K_{d2}.
    $$
    By the same logic applied to the off-target mass-action relation,
    $$
    \frac{Y_C+Y_D}{Y_f}=\frac{C_f+D_f}{K_{d2}}\le\frac{C+D}{K_{d2}},
    $$
    so $(C+D)/K_{d2}$ bounds the ratio of off-target-bound to free
    antibody. Because $K_{d2}\gg K_{d1}$, this assumption is usually much
    easier to satisfy than assumption 1: the off-target pool may be far
    larger than the target pool and still bind little antibody.
\end{enumerate}
In summary, the two assumptions say that the two places antibody can go
other than solution---target and off-target epitopes---each hold a small
amount of antibody compared with the amount that stays free. Note that
neither assumption involves $Y$: in this regime, the antibody may be
limiting, intermediate, or in excess relative to the dissociation
constants; it is only required to be in excess of what the epitopes can
soak up. Because $S_t<S_c$ when $A'<A$, the control condition is the more
demanding one, so it suffices to check the assumptions with $S=S_c$.

We now derive $Y_f$, $\theta_1$, and $\rho$ for either condition. (For the
treatment condition, $\theta_1$, $\theta_2$, $Y_f$, and $\rho$ should be
read as their primed counterparts.)

\textbf{Antibody conservation.} As derived in Regime 1, substituting
$Y_A+Y_B=S\theta_1$ and $Y_C+Y_D=(C+D)\theta_2$ into
$Y=Y_f+Y_A+Y_B+Y_C+Y_D$ gives
\begin{equation}
Y=Y_f+S\theta_1+(C+D)\theta_2.
\label{eq:r2-1}
\end{equation}

\textbf{Step 1: bound the antibody on epitopes.} Dropping $Y_f$ from the
denominator of $\theta_1=Y_f/(Y_f+K_{d1})$ makes the denominator smaller
and hence the fraction larger:
$$
\theta_1=\frac{Y_f}{Y_f+K_{d1}}\le\frac{Y_f}{K_{d1}}.
$$
Multiplying both sides by $S$,
$$
S\theta_1\le S\frac{Y_f}{K_{d1}}=\frac{S}{K_{d1}}\,Y_f.
$$
By the same logic with $\theta_2=Y_f/(Y_f+K_{d2})\le Y_f/K_{d2}$,
$$
(C+D)\theta_2\le\frac{C+D}{K_{d2}}\,Y_f.
$$

\textbf{Step 2: $Y_f\approx Y$.} Define the relative amount of
epitope-bound antibody
$$
\delta=\frac{S\theta_1+(C+D)\theta_2}{Y_f},
$$
so that \eqref{eq:r2-1} reads $Y=Y_f+\delta Y_f=Y_f(1+\delta)$. Adding the
two inequalities from Step 1 and dividing by $Y_f$,
\begin{equation}
0\le\delta
\le\frac{S}{K_{d1}}+\frac{C+D}{K_{d2}}
\equiv\bar\delta
\ll1,
\label{eq:r2-2}
\end{equation}
where $\delta\ge0$ because every term is nonnegative, and the final
$\ll$ follows from assumptions 1 and 2 (a sum of two small numbers is
small). Solving $Y=Y_f(1+\delta)$ for $Y_f$ and using
$1\le1+\delta\le1+\bar\delta$,
\begin{equation}
\frac{Y}{1+\bar\delta}\le Y_f=\frac{Y}{1+\delta}\le Y.
\label{eq:r2-3}
\end{equation}
Since $\bar\delta\ll1$, $Y_f\approx Y$: the free antibody is essentially
all of the antibody. (Equivalently, the fraction of antibody that is
bound, $(Y-Y_f)/Y=\delta/(1+\delta)$, is at most $\bar\delta$.)

\textbf{Step 3: capture efficiencies.} Substituting $Y_f\approx Y$ into
the expressions for the bound fractions,
$$
\theta_1\approx\frac{Y}{Y+K_{d1}},
\qquad
\theta_2\approx\frac{Y}{Y+K_{d2}}.
$$
Neither depends on $S$. In particular, the fraction of target epitope
captured is the same in both conditions, so the IP yield \emph{per
epitope} does not change when the treatment removes target. (The total
target captured, $S\theta_1$, is proportional to $S$ and so does
decrease.) This is the key difference from Regime 1, where the total
captured was fixed at about $Y$ and the per-epitope efficiency
$\theta_1\approx Y/S$ rose when $S$ fell.

\textbf{Step 4: $\rho$.} Substituting $Y_f\approx Y$ into
$\rho=(Y_f+K_{d1})/(Y_f+K_{d2})$ gives
\begin{equation}
\rho\approx\rho'\approx\bar\rho\equiv\frac{Y+K_{d1}}{Y+K_{d2}}.
\label{eq:r2-4}
\end{equation}
Since $\bar\rho$ depends only on $Y$, $K_{d1}$, and $K_{d2}$---not on
$S$---it is the same for the control and treatment conditions,
regardless of $A'$.

To make ``$\approx$'' precise, we bound $\rho$ above and below. First,
$\rho$ increases with $Y_f$: by the quotient rule,
$$
\frac{d\rho}{dY_f}
=\frac{(Y_f+K_{d2})-(Y_f+K_{d1})}{(Y_f+K_{d2})^2}
=\frac{K_{d2}-K_{d1}}{(Y_f+K_{d2})^2}
>0,
$$
since $K_{d2}>K_{d1}$. Because $Y_f\le Y$ by \eqref{eq:r2-3}, it follows
that $\rho\le\bar\rho$. For the lower bound, substitute the lower bound
$Y_f\ge Y/(1+\bar\delta)$ from \eqref{eq:r2-3} and multiply the numerator
and denominator by $1+\bar\delta$:
$$
\rho
\ge\frac{\dfrac{Y}{1+\bar\delta}+K_{d1}}{\dfrac{Y}{1+\bar\delta}+K_{d2}}
=\frac{Y+(1+\bar\delta)K_{d1}}{Y+(1+\bar\delta)K_{d2}}.
$$
The numerator is at least $Y+K_{d1}$, and the denominator is at most
$(1+\bar\delta)Y+(1+\bar\delta)K_{d2}=(1+\bar\delta)(Y+K_{d2})$, so
$$
\rho\ge\frac{Y+K_{d1}}{(1+\bar\delta)(Y+K_{d2})}=\frac{\bar\rho}{1+\bar\delta}.
$$
Together,
\begin{equation}
\frac{\bar\rho}{1+\bar\delta}\le\rho\le\bar\rho,
\label{eq:r2-5}
\end{equation}
so $\rho$ agrees with $\bar\rho$ to within a relative error of at most
$\bar\delta\ll1$. The same holds for $\rho'$ (with $\bar\delta$ computed
using $S_t$, which is smaller than with $S_c$).

\textbf{Direction of the (small) difference.} Although $\rho'\approx\rho$,
the exact values still move in the same direction as in Regime 1. Write
the right-hand side of \eqref{eq:r2-1} as a function of $Y_f$ and $S$:
$$
g(Y_f,S)=Y_f+S\frac{Y_f}{Y_f+K_{d1}}+(C+D)\frac{Y_f}{Y_f+K_{d2}}.
$$
Each term increases with $Y_f$ (for instance,
$\frac{d}{dY_f}\frac{Y_f}{Y_f+K}=\frac{K}{(Y_f+K)^2}>0$), and $g$ increases
with $S$. If $A'<A$, then $S_t<S_c$, so
$$
g(Y_f,S_t)<g(Y_f,S_c)=Y=g(Y_f',S_t).
$$
Since $g(\cdot,S_t)$ is increasing, this requires $Y_f<Y_f'$: with less
target to bind, slightly more antibody stays free. Because $\rho$
increases with $Y_f$, $\rho<\rho'$. Combining with \eqref{eq:r2-5} for
both conditions,
\begin{equation}
\rho<\rho'\le\bar\rho\le(1+\bar\delta)\rho,
\label{eq:r2-6}
\end{equation}
where $\bar\delta$ is evaluated for the control ($S=S_c$). Thus $\rho'$
exceeds $\rho$, but by a relative amount of at most $\bar\delta\ll1$.

\textbf{Sub-limits of $\bar\rho$.} Factoring $K_{d1}$ out of the numerator
and $K_{d2}$ out of the denominator of \eqref{eq:r2-4},
$$
\bar\rho
=\frac{K_{d1}(1+Y/K_{d1})}{K_{d2}(1+Y/K_{d2})}
=r\,\frac{1+Y/K_{d1}}{1+Y/K_{d2}}.
$$
\begin{itemize}
    \item $Y\ll K_{d1}$ (antibody below the target $K_d$): both
    $Y/K_{d1}\ll1$ and $Y/K_{d2}\ll1$, so $\bar\rho\approx r$. Both
    epitope classes are unsaturated, $\theta_1\approx Y/K_{d1}$ and
    $\theta_2\approx Y/K_{d2}$, and their ratio is set purely by the
    affinities.
    \item $K_{d1}\ll Y\ll K_{d2}$: in the numerator,
    $Y+K_{d1}\approx Y$; in the denominator, $Y+K_{d2}\approx K_{d2}$. So
    $\bar\rho\approx Y/K_{d2}$. The target is saturated
    ($\theta_1\approx1$), while off-target binding is still in its linear
    regime ($\theta_2\approx Y/K_{d2}$), so off-target capture grows in
    proportion to $Y$ while target capture cannot grow further.
    \item $Y\gg K_{d2}$: both $Y+K_{d1}\approx Y$ and $Y+K_{d2}\approx Y$,
    so $\bar\rho\approx1$. Both classes are saturated and the antibody no
    longer discriminates between them.
\end{itemize}
As $Y$ increases, $\bar\rho$ rises monotonically from $r$ to 1 (by the
same derivative computed above, with $Y$ in place of $Y_f$).

\subsubsection{ChIP normalization}

The read-count expressions derived in Regime 1 (expected reads $n_i$ at
target locus $i$ and $N_s$ on the spike-in genome) did not use any
regime-specific assumptions, so they apply here unchanged.

\textbf{Scaling by spike-in reads (ChIP-Rx).} From Regime 1,
$$
\frac{n_i'/N_s'}{n_i/N_s}
=\frac{a_i'}{a_i}\cdot\frac{B+\rho D}{B+\rho' D}.
$$
Substituting $\rho\approx\rho'\approx\bar\rho$ from \eqref{eq:r2-4}, the
numerator and denominator of the bias factor are both approximately
$B+\bar\rho D$, so
$$
\frac{B+\rho D}{B+\rho' D}\approx\frac{B+\bar\rho D}{B+\bar\rho D}=1,
\qquad
\frac{n_i'/N_s'}{n_i/N_s}\approx\frac{a_i'}{a_i}.
$$
More precisely, by \eqref{eq:r2-6}, $\rho D<\rho' D\le(1+\bar\delta)\rho D$.
Adding $B$ and using $B\le(1+\bar\delta)B$,
$$
B+\rho D<B+\rho' D\le(1+\bar\delta)(B+\rho D),
$$
and taking reciprocals and multiplying by $B+\rho D$,
$$
\frac{1}{1+\bar\delta}\le\frac{B+\rho D}{B+\rho' D}<1.
$$
So spike-in normalization recovers the true fold change to within a
relative error of at most $\bar\delta$, \emph{even if off-target spike-in
reads are not negligible} ($\rho D$ comparable to or larger than $B$).
This contrasts with Regime 1, where the bias factor equals 1 only if
$\rho D\ll B$. The reason is that in this regime the off-target spike-in
reads are captured with the same efficiency in both conditions, so they
scale the spike-in read count by the same factor in both libraries and
cancel in the ratio. (The direction of the residual bias, $<1$, is the
same as in Regime 1, but its size is bounded by $\bar\delta$.) Scaling by
spike-in target reads $N_{s,B}$ also works, as it did in Regime 1.

\textbf{Scaling by total reads (RPM).} From Regime 1, substituting
$\rho'\approx\rho$,
$$
\frac{n_i'/N'}{n_i/N}
=\frac{a_i'}{a_i}\cdot\frac{A+B+\rho(C+D)}{A'+B+\rho'(C+D)}
\approx\frac{a_i'}{a_i}\cdot\frac{A+B+\rho(C+D)}{A'+B+\rho(C+D)}.
$$
If $A'<A$, the numerator of the second factor exceeds its denominator, so
the factor is greater than 1 and partially cancels the global decrease
(fully, $\approx A/A'$, if the spike-in and off-target terms are small).
RPM normalization still hides global changes in this regime.

\textbf{The cost of this regime.} Dilute epitopes make normalization
robust, but they come at a price.
\begin{itemize}
    \item \emph{Low yield when $Y\ll K_{d1}$.} The fraction of target
    captured is $\theta_1\approx Y/(Y+K_{d1})\approx Y/K_{d1}\ll1$, and the
    total target captured is $S\theta_1\approx SY/K_{d1}$, which is small
    because both $S\ll K_{d1}$ and $Y\ll K_{d1}$. Little DNA is recovered,
    so libraries have low complexity and are sensitive to contamination.
    \item \emph{Off-target contamination when $Y$ is raised.} Increasing $Y$
    to improve yield raises $\bar\rho$ from $r$ toward $Y/K_{d2}$ and
    eventually 1. Among endogenous reads, the target-to-off-target ratio is
    $A:\bar\rho C$, so the fraction of endogenous reads that are on target,
    $$
    \frac{A}{A+\bar\rho C},
    $$
    falls as $Y$ grows. Once $Y\gg K_{d1}$, target capture is saturated
    ($\theta_1\approx1$), so further increases in $Y$ only add off-target
    background. Because $\bar\rho$ is the same in both conditions, this
    background does not bias spike-in normalization at target loci (under
    the assumption that target and off-target loci do not overlap), but it
    lowers the signal-to-noise ratio, especially in the treatment
    condition, where the on-target fraction $A'/(A'+\bar\rho C)$ is
    smaller because $A'<A$.
\end{itemize}

\subsection{Regime 3: antibody in excess of target}

As in Regime 1, let $S_c=A+B$ and $S_t=A'+B$ be the total target-epitope
concentrations in the control and treatment conditions, and let
$S\in\{S_c,S_t\}$ denote either one. We assume the following for both
conditions (i.e., for each $S\in\{S_c,S_t\}$).

The assumptions are again easiest to interpret through a heuristic
estimate of the free-antibody concentration. Now there is more antibody
than target. If nearly every target epitope is bound, then about $S$
antibody molecules sit on target, and the excess $Y-S$ must go somewhere
else: either into solution ($Y_f$) or onto off-target epitopes
($Y_C+Y_D$). If off-target epitopes are far from saturated, nearly all of
them are free ($C_f+D_f\approx C+D$), so rearranging the mass-action
relation $K_{d2}=Y_f(C_f+D_f)/(Y_C+Y_D)$ gives
$$
Y_C+Y_D
=\frac{Y_f(C_f+D_f)}{K_{d2}}
\approx\frac{Y_f(C+D)}{K_{d2}}.
$$
Substituting this into ``excess $=$ free $+$ off-target-bound,''
$$
Y-S
\approx Y_f+\frac{Y_f(C+D)}{K_{d2}}
=Y_f\cdot\frac{K_{d2}+C+D}{K_{d2}},
$$
and multiplying both sides by $K_{d2}/(K_{d2}+C+D)$,
$$
Y_f\approx K_{d2}\frac{Y-S}{K_{d2}+C+D}.
$$
In words, the excess antibody is split between solution and off-target
epitopes in the ratio $K_{d2}:(C+D)$. Below, we show that this estimate
is in fact an exact lower bound on $Y_f$.

\begin{enumerate}
    \item Antibody is more abundant than target epitope:
    $$
    Y>S.
    $$
    Thus, even if every target epitope is bound, there is a positive
    excess $Y-S$ of antibody left over.

    \item Free target is scarce relative to the excess antibody:
    $$
    rS\,(K_{d2}+C+D)\ll(Y-S)^2.
    $$
    To interpret this, estimate the free target concentration
    $A_f+B_f$. Rearranging the mass-action relation
    $K_{d1}=Y_f(A_f+B_f)/(Y_A+Y_B)$, and using that about $S$ target
    epitopes are bound,
    $$
    A_f+B_f
    =K_{d1}\frac{Y_A+Y_B}{Y_f}
    \approx\frac{K_{d1}S}{Y_f}.
    $$
    Substituting the heuristic estimate of $Y_f$ and then
    $r=K_{d1}/K_{d2}$,
    $$
    A_f+B_f
    \approx\frac{K_{d1}S\,(K_{d2}+C+D)}{K_{d2}(Y-S)}
    =\frac{rS\,(K_{d2}+C+D)}{Y-S}.
    $$
    Dividing both sides of the condition by $Y-S$ shows that it is
    equivalent to
    $rS(K_{d2}+C+D)/(Y-S)\ll Y-S$: the estimated free target is much less
    than the excess antibody. This matters because every target epitope
    that is free leaves one more antibody molecule unaccounted for, so the
    antibody available to go into solution or off-target stays close to
    $Y-S$ only if the free target is small compared with $Y-S$. Usually
    this also means that the target is nearly saturated
    ($\theta_1\approx1$). The square arises because $Y-S$ appears twice:
    once in the estimate of the free target (more excess antibody drives
    more target into the bound state) and once as the yardstick against
    which the free target is compared.

    \textbf{Mirror symmetry with Regime 1.} Substituting
    $r=K_{d1}/K_{d2}$ and splitting the fraction, the left-hand side is
    $$
    rS\,(K_{d2}+C+D)
    =\frac{K_{d1}}{K_{d2}}S\,(K_{d2}+C+D)
    =K_{d1}S\left(1+\frac{C+D}{K_{d2}}\right).
    $$
    If there were no off-target epitopes ($C+D=0$), the condition would
    read $K_{d1}S\ll(Y-S)^2$. Compare Regime 1's assumption 2,
    $K_{d1}Y\ll(S-Y)^2$: the two are identical after swapping $Y$ and
    $S$. This is because the mass-action relation
    $K_{d1}=Y_f(A_f+B_f)/(Y_A+Y_B)$ treats antibody and target
    symmetrically. In Regime 1, antibody is the limiting species and target
    is in excess; in Regime 3, the roles are reversed. In both cases the
    condition says that the leftover \emph{limiting} species (free antibody
    in Regime 1, free target here) is much less than the excess of the
    \emph{abundant} species. The extra factor $1+(C+D)/K_{d2}$ breaks the
    symmetry: off-target epitopes are a sink for antibody but not for
    target, so only a fraction $K_{d2}/(K_{d2}+C+D)$ of the excess
    antibody remains free to drive target binding. A weaker push on the
    target requires a larger excess to compensate.

    Formally, below we show that $S(1-\theta_1)\le rS(K_{d2}+C+D)/(Y-S)$,
    so this condition makes the free target negligible relative to $Y-S$.

    \item Off-target epitopes are far from saturated:
    $$
    Y-S\ll K_{d2}+C+D.
    $$
    By the heuristic estimate, the off-target occupancy is
    $$
    \theta_2=\frac{Y_f}{Y_f+K_{d2}}\approx\frac{Y_f}{K_{d2}}
    \approx\frac{Y-S}{K_{d2}+C+D},
    $$
    so this condition says that the fraction of off-target epitopes that
    are bound is small. Another way to see the yardstick
    $K_{d2}+C+D$: to half-occupy the off-target epitopes, the free antibody
    must reach $K_{d2}$ and about $(C+D)/2$ antibody molecules must be bound
    off-target, for a total of about $K_{d2}+(C+D)/2$, which is within a
    factor of 2 of $K_{d2}+C+D$. The condition says that the excess
    antibody is far below this amount. Below, we show that this condition
    implies $Y_f\ll K_{d2}$.
\end{enumerate}
In summary, assumption 2 ensures that essentially all of the excess
antibody $Y-S$ is available to the ``non-target'' compartments (solution
and off-target epitopes), and assumption 3 keeps the off-target epitopes
in their unsaturated, linear-response regime. Compared with Regime 1,
there is no separate assumption limiting off-target binding: here,
off-target binding is not a perturbation to be neglected but one of the
two places where the excess antibody goes.

\textbf{Simpler sufficient condition.} Suppose $Y\ge2S$. Subtracting $S$
from both sides gives $Y-S\ge S$, and since $Y-S>0$ by assumption 1,
dividing both sides by $Y-S$ gives
$$
\frac{S}{Y-S}\le1.
$$
Writing the left-hand side of assumption 2, divided by $(Y-S)^2$, as a
product,
$$
\frac{rS\,(K_{d2}+C+D)}{(Y-S)^2}
=\frac{r(K_{d2}+C+D)}{Y-S}\cdot\frac{S}{Y-S}
\le\frac{r(K_{d2}+C+D)}{Y-S}.
$$
Therefore, if $r(K_{d2}+C+D)\ll Y-S$, then assumption 2 holds. When $Y$
approaches $S$, however, the factor $S/(Y-S)$ grows without bound, and
the simpler condition is no longer sufficient.

We now derive $\rho$ for either condition. (For the treatment condition,
$\theta_1$, $\theta_2$, $Y_f$, and $\rho$ should be read as their primed
counterparts.)

\textbf{Antibody conservation.} As in Regime 1,
$$
Y=Y_f+S\theta_1+(C+D)\theta_2.
$$

\textbf{Step 1: $\rho$ in terms of $\theta_2$.} Substituting
$K_{d1}=rK_{d2}$ into the numerator of $\rho$, and then adding and
subtracting $rY_f$,
$$\begin{aligned}
\rho
&=\frac{Y_f+K_{d1}}{Y_f+K_{d2}}
=\frac{Y_f+rK_{d2}}{Y_f+K_{d2}} \\
&=\frac{r(Y_f+K_{d2})+Y_f-rY_f}{Y_f+K_{d2}} \\
&=r+(1-r)\frac{Y_f}{Y_f+K_{d2}}.
\end{aligned}$$
Recognizing $\theta_2=Y_f/(Y_f+K_{d2})$,
\begin{equation}
\rho=r+(1-r)\,\theta_2.
\label{eq:r3-1}
\end{equation}
This is exact and holds in every regime. It says that $\rho$ interpolates
linearly between $r$ (no off-target occupancy) and 1 (full off-target
occupancy). So it remains to find $\theta_2$.

\textbf{Step 2: where the excess antibody goes.} Write the
target-bound antibody as total target minus free target,
$S\theta_1=S-S(1-\theta_1)$. Substituting into antibody conservation,
$$
Y=Y_f+S-S(1-\theta_1)+(C+D)\theta_2.
$$
Subtracting $S$ from both sides and adding $S(1-\theta_1)$ to both sides,
\begin{equation}
(Y-S)+S(1-\theta_1)=Y_f+(C+D)\theta_2.
\label{eq:r3-2}
\end{equation}
Intuitively, the antibody not on target (right-hand side) equals the
excess antibody $Y-S$ plus one antibody for each target epitope that is
left free. This is the mirror image of Regime 1's equation
\eqref{eq:r1-1}.

\textbf{Step 3: an exact expression for $\theta_2$.} From
$\theta_2=Y_f/(Y_f+K_{d2})$, multiplying both sides by $Y_f+K_{d2}$ gives
$Y_f=\theta_2(Y_f+K_{d2})$. Substituting this into the right-hand side of
\eqref{eq:r3-2} and factoring out $\theta_2$,
$$
(Y-S)+S(1-\theta_1)
=\theta_2(Y_f+K_{d2})+(C+D)\theta_2
=\theta_2\,(K_{d2}+C+D+Y_f).
$$
Dividing both sides by $K_{d2}+C+D+Y_f$,
$$
\theta_2=\frac{(Y-S)+S(1-\theta_1)}{K_{d2}+C+D+Y_f}.
$$
Factoring $Y-S$ out of the numerator and $K_{d2}+C+D$ out of the
denominator,
\begin{equation}
\theta_2
=\frac{Y-S}{K_{d2}+C+D}\cdot\frac{1+\delta_1}{1+\delta_2},
\qquad
\delta_1=\frac{S(1-\theta_1)}{Y-S},
\quad
\delta_2=\frac{Y_f}{K_{d2}+C+D}.
\label{eq:r3-3}
\end{equation}
Equations \eqref{eq:r3-1} and \eqref{eq:r3-3} together are exact. Here
$\delta_1$ is the free target relative to the excess antibody, and
$\delta_2$ is the free antibody relative to $K_{d2}+C+D$. It remains to
show that $\delta_1\ll1$ and $\delta_2\ll1$.

\textbf{Step 4: a lower bound on $Y_f$.} Because $S(1-\theta_1)\ge0$,
equation \eqref{eq:r3-2} gives
$$
Y_f+(C+D)\theta_2\ge Y-S.
$$
Dropping $Y_f$ from the denominator of $\theta_2$ makes it larger, so
$\theta_2=Y_f/(Y_f+K_{d2})\le Y_f/K_{d2}$. Substituting this on the
left-hand side (which can only make the left-hand side larger, so the
inequality still holds),
$$
Y_f+\frac{(C+D)Y_f}{K_{d2}}\ge Y-S
\quad\Longrightarrow\quad
Y_f\cdot\frac{K_{d2}+C+D}{K_{d2}}\ge Y-S.
$$
Multiplying both sides by the positive number $K_{d2}/(K_{d2}+C+D)$,
\begin{equation}
Y_f\ge K_{d2}\frac{Y-S}{K_{d2}+C+D},
\label{eq:r3-4}
\end{equation}
which is positive by assumption 1. This is the heuristic estimate from
the start of this subsection, now shown to be a lower bound.

\textbf{Step 5: $\delta_1\ll1$.} From $\theta_1=Y_f/(Y_f+K_{d1})$, as in
Regime 1, $1-\theta_1=K_{d1}/(Y_f+K_{d1})$. Dropping $K_{d1}$ from the
denominator makes the fraction larger, and then substituting the lower
bound \eqref{eq:r3-4} for $Y_f$ in the denominator makes it larger
still:
$$
1-\theta_1
=\frac{K_{d1}}{Y_f+K_{d1}}
\le\frac{K_{d1}}{Y_f}
\le\frac{K_{d1}(K_{d2}+C+D)}{K_{d2}(Y-S)}
=\frac{r(K_{d2}+C+D)}{Y-S},
$$
where the last step substitutes $r=K_{d1}/K_{d2}$. Multiplying both sides
by $S/(Y-S)$,
$$
\delta_1
=\frac{S(1-\theta_1)}{Y-S}
\le\frac{rS\,(K_{d2}+C+D)}{(Y-S)^2}
\ll1,
$$
by assumption 2.

\textbf{Step 6: $\theta_2\ll1$, $Y_f\ll K_{d2}$, and $\delta_2\ll1$.}
Because $\delta_2\ge0$, the factor $1/(1+\delta_2)$ in \eqref{eq:r3-3} is
at most 1, so
$$
\theta_2\le\frac{Y-S}{K_{d2}+C+D}\,(1+\delta_1)\ll1,
$$
since $(Y-S)/(K_{d2}+C+D)\ll1$ by assumption 3 and $1+\delta_1\approx1$ by
Step 5. To convert this into a statement about $Y_f$, start from
$Y_f=\theta_2(Y_f+K_{d2})$, subtract $\theta_2Y_f$ from both sides to get
$Y_f(1-\theta_2)=\theta_2K_{d2}$, and divide by $K_{d2}(1-\theta_2)$:
$$
\frac{Y_f}{K_{d2}}=\frac{\theta_2}{1-\theta_2}.
$$
With $\theta_2\ll1$, the right-hand side is approximately $\theta_2\ll1$,
so $Y_f\ll K_{d2}$. Finally, because $K_{d2}\le K_{d2}+C+D$,
$$
\delta_2=\frac{Y_f}{K_{d2}+C+D}\le\frac{Y_f}{K_{d2}}\ll1.
$$

\textbf{Step 7: result.} Since $\delta_1\ll1$ and $\delta_2\ll1$,
$(1+\delta_1)/(1+\delta_2)\approx1$, and \eqref{eq:r3-3} becomes
$$
\theta_2\approx\frac{Y-S}{K_{d2}+C+D}.
$$
Substituting this into \eqref{eq:r3-1}, and using $1-r\approx1$ (because
$K_{d1}\ll K_{d2}$),
\begin{equation}
\rho\approx r+\frac{Y-S}{K_{d2}+C+D}.
\label{eq:r3-5}
\end{equation}
Correspondingly, since $Y_f/K_{d2}=\theta_2/(1-\theta_2)\approx\theta_2$,
the free antibody is $Y_f\approx K_{d2}(Y-S)/(K_{d2}+C+D)$, confirming the
heuristic estimate. Substituting $S=A+B$ for the control and $S=A'+B$
for the treatment,
$$
\rho\approx r+\frac{Y-A-B}{K_{d2}+C+D},
\qquad
\rho'\approx r+\frac{Y-A'-B}{K_{d2}+C+D}.
$$

(An exact lower bound also follows. Since
$d\theta_2/dY_f=K_{d2}/(Y_f+K_{d2})^2>0$, $\theta_2$ increases with
$Y_f$, so substituting the lower bound \eqref{eq:r3-4} into
$\theta_2=Y_f/(Y_f+K_{d2})$ and then multiplying the numerator and
denominator by $(K_{d2}+C+D)/K_{d2}$,
$$
\theta_2
\ge\frac{K_{d2}\frac{Y-S}{K_{d2}+C+D}}
{K_{d2}\frac{Y-S}{K_{d2}+C+D}+K_{d2}}
=\frac{Y-S}{(Y-S)+(K_{d2}+C+D)}.
$$
Substituting into \eqref{eq:r3-1} gives
$\rho\ge r+(1-r)(Y-S)/(Y-S+K_{d2}+C+D)$, which agrees with
\eqref{eq:r3-5} to first order when $Y-S\ll K_{d2}+C+D$.)

Two terms contribute to $\rho$. The first, $r$, is the baseline ratio of
off-target to target capture efficiency set by the affinities. The
second is the fraction of off-target epitopes occupied by the excess
antibody. Because target is saturated ($\theta_1\approx1$), any increase
in off-target occupancy translates directly into an increase in $\rho$.
Unlike in Regime 1, the second term is not multiplied by $r$, so $\rho$ is
not necessarily close to $r$: it can be many times larger.

\textbf{Two sub-cases.} The denominator $K_{d2}+C+D$ reflects the two
places where the excess antibody can go, in the ratio $K_{d2}:(C+D)$.
\begin{itemize}
    \item[(3a)] Few off-target epitopes, $C+D\ll K_{d2}$. Then
    $K_{d2}+C+D\approx K_{d2}$, so
    $$
    Y_f\approx K_{d2}\frac{Y-S}{K_{d2}}=Y-S,
    \qquad
    \rho\approx r+\frac{Y-S}{K_{d2}}.
    $$
    The excess antibody stays free in solution, and the off-target
    occupancy is simply $Y_f/K_{d2}$. Assumption 3 becomes $Y-S\ll K_{d2}$,
    and assumption 2 becomes $K_{d1}S\ll(Y-S)^2$, the exact mirror of
    Regime 1's assumption 2.

    \item[(3b)] Many off-target epitopes, $C+D\gg K_{d2}$. Then
    $K_{d2}+C+D\approx C+D$, so
    $$
    Y_f\approx K_{d2}\frac{Y-S}{C+D}\ll Y-S,
    \qquad
    \rho\approx r+\frac{Y-S}{C+D}.
    $$
    The off-target pool soaks up nearly all of the excess antibody, so the
    excess divided by the number of off-target epitopes is the fraction of
    off-target epitopes that are occupied. Assumption 3 becomes
    $Y-S\ll C+D$.
\end{itemize}
Chromatin contains an enormous number of weak, nonspecific binding sites,
so sub-case (3b) is plausible in practice.

(Remark on assumption 3. The approximation to $\theta_2$ in
\eqref{eq:r3-3} actually requires only $\delta_2\ll1$, i.e.,
$Y_f\ll K_{d2}+C+D$. In sub-case (3b), this can remain true even when a
substantial fraction of off-target epitopes is occupied, because the
large off-target pool keeps $Y_f$ small. Then the exact relation
\eqref{eq:r3-1} with $\theta_2\approx(Y-S)/(C+D)$ still holds, and $\rho$
continues to grow linearly with the excess until $Y-S$ approaches $C+D$.
We keep assumption 3 as stated because it defines the regime in which
$\rho\ll1$.)

\textbf{How $\rho'$ compares to $\rho$.} Subtracting the approximations
for $\rho$ and $\rho'$, the $r$ terms cancel and the fractions share a
common denominator:
$$\begin{aligned}
\rho'-\rho
&\approx\left(r+\frac{Y-S_t}{K_{d2}+C+D}\right)
-\left(r+\frac{Y-S_c}{K_{d2}+C+D}\right) \\
&=\frac{(Y-S_t)-(Y-S_c)}{K_{d2}+C+D}
=\frac{S_c-S_t}{K_{d2}+C+D}.
\end{aligned}$$
Substituting $S_c=A+B$ and $S_t=A'+B$, the spike-in terms $B$ cancel:
\begin{equation}
\rho'-\rho\approx\frac{(A+B)-(A'+B)}{K_{d2}+C+D}
=\frac{A-A'}{K_{d2}+C+D}.
\label{eq:r3-6}
\end{equation}
Since the denominator is positive, \textbf{$\rho'>\rho$ when $A'<A$}.
Intuitively, the treatment removes $A-A'$ target epitopes, which frees
the $A-A'$ antibody molecules that would have occupied them. That
antibody joins the excess, and a fraction $(C+D)/(K_{d2}+C+D)$ of it
lands on off-target epitopes, raising the off-target occupancy.

The sign of $\rho'-\rho$ is in fact exact. Define
$g(Y_f)=Y_f+SY_f/(Y_f+K_{d1})+(C+D)Y_f/(Y_f+K_{d2})$, so that antibody
conservation reads $g(Y_f)=Y$. Each term of $g$ increases with $Y_f$,
and for $Y_f>0$ the middle term increases with $S$. If $S$ decreases
(from $S_c$ to $S_t$) while $Y$ is held fixed, then $g$ evaluated at the
old $Y_f$ falls below $Y$, so the new solution must have a larger $Y_f$.
By the quotient rule,
$$
\frac{d\rho}{dY_f}
=\frac{(Y_f+K_{d2})-(Y_f+K_{d1})}{(Y_f+K_{d2})^2}
=\frac{K_{d2}-K_{d1}}{(Y_f+K_{d2})^2}>0,
$$
so a larger $Y_f$ means a larger $\rho$. This argument applies in every
regime, including Regime 1.

Note that the magnitude \eqref{eq:r3-6}, unlike the Regime 1 result, is
not suppressed by a factor of $r$. In Regime 1, a decrease in target
changes $\rho$ only by changing how tightly the (limiting) antibody is
held on target, an effect proportional to $r$. Here, the decrease in
target releases antibody directly into the off-target pool.

\textbf{Relevance to practice.} Typical ChIP protocols use micrograms of
antibody. An IgG has a molar mass of about 150~kDa, so
$1~\mu\mathrm{g}$ is about
$10^{-6}~\mathrm{g}/(1.5\times10^{5}~\mathrm{g/mol})\approx6.7~\mathrm{pmol}$.
By comparison, a factor with $10^4$ binding sites per cell in $10^6$ cells
has $10^{10}$ sites, or about
$10^{10}/(6.0\times10^{23}~\mathrm{mol}^{-1})\approx0.017~\mathrm{pmol}$.
For many targets, then, standard practice places the IP in this
antibody-excess regime. (Very abundant epitopes, such as some histone
modifications, may instead be comparable to or exceed the antibody.)

\subsubsection{ChIP normalization}

The read-count model of Regime 1 applies unchanged, because it used only
the composition $A:B:\rho C:\rho D$ of the immunoprecipitation. In
particular, ChIP-Rx scaling gives
$$
\frac{n_i'/N_s'}{n_i/N_s}
=\frac{a_i'}{a_i}\cdot\frac{B+\rho D}{B+\rho' D}.
$$

\textbf{Rewriting the bias factor.} Write $\rho'=\rho+(\rho'-\rho)$, so
the denominator is
$B+\rho'D=(B+\rho D)+(\rho'-\rho)D$. Dividing the numerator and
denominator of the bias factor by $B+\rho D$,
$$
\frac{B+\rho D}{B+\rho'D}
=\frac{B+\rho D}{(B+\rho D)+(\rho'-\rho)D}
=\frac{1}{1+\dfrac{(\rho'-\rho)D}{B+\rho D}}.
$$
Since $\rho'>\rho$ when $A'<A$, the fraction in the denominator is
positive, so the bias factor is less than 1: as in Regime 1,
\textbf{ChIP-Rx overstates the decrease} in the treatment condition. The
bias factor is approximately 1 only if
$$
(\rho'-\rho)D\ll B+\rho D.
$$

\textbf{Substituting the regime expressions.} Substituting
\eqref{eq:r3-6} for $\rho'-\rho$ and \eqref{eq:r3-5} for $\rho$, the
bias factor is
$$
\frac{B+\rho D}{B+\rho'D}
\approx\frac{1}{1+\dfrac{(A-A')D}{(K_{d2}+C+D)\bigl(B+\rho D\bigr)}},
\qquad
\rho\approx r+\frac{Y-A-B}{K_{d2}+C+D},
$$
and ChIP-Rx is unbiased only if
$$
\frac{(A-A')D}{K_{d2}+C+D}\ll B+\left(r+\frac{Y-A-B}{K_{d2}+C+D}\right)D.
$$
The left-hand side has a direct interpretation. Splitting the fraction,
$$
\frac{(A-A')D}{K_{d2}+C+D}
=(A-A')\cdot\frac{C+D}{K_{d2}+C+D}\cdot\frac{D}{C+D}:
$$
the $A-A'$ antibody molecules freed by the treatment, times the fraction
of excess antibody that is captured off-target, times the share of
off-target epitopes that belong to the spike-in. That is the extra
antibody that ends up on spike-in off-target epitopes in the treatment.
The right-hand side, $B+\rho D=(B\theta_1+D\theta_2)/\theta_1$, is the
spike-in material in the control IP divided by $\theta_1\approx1$. So the
condition says that the extra spike-in capture caused by the freed
antibody must be small compared with the spike-in capture already
present.

\textbf{How large can the bias be?} Consider sub-case (3b), so
$K_{d2}+C+D\approx C+D$, and suppose off-target spike-in reads are a
minor part of the spike-in signal, $\rho D\ll B$. The fraction in the
bias factor becomes
$$
\frac{(A-A')D}{(K_{d2}+C+D)(B+\rho D)}
\approx\frac{(A-A')D}{(C+D)B}.
$$
Suppose further that the spike-in chromatin has the same ratio of
off-target to target epitopes as the cells, $D/B=C/A$. Substituting
$D/B=C/A$,
$$
\frac{(A-A')D}{(C+D)B}
=\frac{(A-A')}{C+D}\cdot\frac{C}{A}
=\frac{A-A'}{A}\cdot\frac{C}{C+D}.
$$
If the spike-in is a small fraction of the chromatin ($D\ll C$), then
$C/(C+D)\approx1$ and the fraction is approximately $1-A'/A$, the
fractional loss of target. The bias factor is then
$$
\frac{1}{1+(1-A'/A)}=\frac{1}{2-A'/A},
$$
and the measured global fold change is
$$
\frac{A'}{A}\cdot\frac{1}{2-A'/A}.
$$
For example, if the treatment halves the target ($A'/A=1/2$), the bias
factor is $1/(2-1/2)=2/3$, and ChIP-Rx reports a fold change of
$(1/2)(2/3)=1/3$ instead of $1/2$. This bias is of order 1, whereas in
Regime 1 it was suppressed by $r$. The reason is that in this regime the
freed antibody is redistributed almost entirely onto off-target
epitopes, and the spike-in receives a share of that redistribution
comparable to its share of the target signal. (If instead $\rho D\gg B$,
the bias factor approaches $\rho D/(\rho' D)=\rho/\rho'$, the ratio of
off-target capture efficiencies, which can also differ substantially
from 1 when the excess $Y-S$ is not large compared with $A-A'$.)

\textbf{Scaling by spike-in target reads.} The derivation of Regime 1
used only the read-count model, not any regime-specific approximation.
Therefore, counting only reads $N_{s,B}$ at the spike-in's target loci
still gives
$$
\frac{n_i'/N_{s,B}'}{n_i/N_{s,B}}=\frac{a_i'}{a_i},
$$
removing the bias regardless of $\rho$ and $\rho'$. This matters more
here than in Regime 1, because the off-target spike-in reads that
$N_s$ includes are exactly what changes between conditions. (As before,
reads in spike-in peaks include some off-target background, so the
cancellation is only approximate; the residual error is larger when
$\rho D$ is larger.)

\subsection{Regime 4: saturating antibody}

In this regime, antibody is so abundant that even the weak, off-target
interactions are saturated. As in Regime 1, let $S\in\{S_c,S_t\}$ with
$S_c=A+B$ and $S_t=A'+B$, and assume the following for both conditions.

\textbf{Assumption.} The antibody left over after binding \emph{every}
epitope, target and off-target, is much greater than $K_{d2}$:
$$
Y-S-(C+D)\gg K_{d2}.
$$
For brevity, write $E=Y-S-(C+D)$ for this \textbf{excess antibody}
($E_c$ and $E_t$ for the two conditions). Because $S_t<S_c$ when $A'<A$,
$$
E_t=Y-S_t-(C+D)>Y-S_c-(C+D)=E_c,
$$
so it suffices to check the assumption for the control condition,
$E_c\gg K_{d2}$.

\textbf{Intuition.} Each epitope can hold at most one antibody, so at
most $S+(C+D)$ antibody can be bound, and at least $E$ must remain free.
Recall that $K_{d2}$ is the free-antibody concentration at which
off-target epitopes are half-bound. If the guaranteed free antibody $E$
is already far above $K_{d2}$, then off-target epitopes are nearly all
bound---and target epitopes, which bind more tightly ($K_{d1}<K_{d2}$),
are even more completely bound. Note that this requires both $Y\gg K_{d2}$
\emph{and} $Y>S+C+D$: antibody must outnumber all epitopes, including the
typically very abundant off-target ones.

We now derive $\rho$ for either condition. (For the treatment condition,
$Y_f$, $\theta_1$, $\theta_2$, $E$, and $\rho$ should be read as their
primed counterparts.)

\textbf{Step 1: a lower bound on $Y_f$.} Start from antibody conservation
(derived in Regime 1),
$$
Y=Y_f+S\theta_1+(C+D)\theta_2.
$$
Solving for $Y_f$, and then adding and subtracting $S+(C+D)$,
$$\begin{aligned}
Y_f
&=Y-S\theta_1-(C+D)\theta_2 \\
&=\bigl[Y-S-(C+D)\bigr]+S-S\theta_1+(C+D)-(C+D)\theta_2 \\
&=E+S(1-\theta_1)+(C+D)(1-\theta_2).
\end{aligned}$$
Intuitively, the free antibody is the excess antibody $E$, plus one
antibody ``returned to solution'' for each epitope that is unbound.
Because $\theta_1\le1$ and $\theta_2\le1$, the last two terms are
nonnegative, so
\begin{equation}
Y_f\ge E.
\label{eq:r4-1}
\end{equation}
By the assumption, $E\gg K_{d2}$, so $Y_f\gg K_{d2}$. Since $K_{d1}<K_{d2}$,
also $Y_f\gg K_{d1}$.

\textbf{Step 2: both classes of epitope are nearly fully bound.} From
$\theta_2=Y_f/(Y_f+K_{d2})$,
$$
1-\theta_2
=\frac{Y_f+K_{d2}}{Y_f+K_{d2}}-\frac{Y_f}{Y_f+K_{d2}}
=\frac{K_{d2}}{Y_f+K_{d2}}.
$$
Dropping $K_{d2}$ from the denominator makes the fraction larger, and
then substituting \eqref{eq:r4-1} (a larger denominator makes it smaller),
$$
1-\theta_2
\le\frac{K_{d2}}{Y_f}
\le\frac{K_{d2}}{E}
\ll1.
$$
By the same logic, $1-\theta_1=K_{d1}/(Y_f+K_{d1})\le K_{d1}/E$, which is
even smaller because $K_{d1}<K_{d2}$. Thus $\theta_1\approx\theta_2\approx1$:
nearly every epitope, target or off-target, is bound by antibody.

\textbf{Step 3: an exact expression for $1-\rho$.} Writing $1$ with the
same denominator as $\rho$,
\begin{equation}
1-\rho
=\frac{Y_f+K_{d2}}{Y_f+K_{d2}}-\frac{Y_f+K_{d1}}{Y_f+K_{d2}}
=\frac{K_{d2}-K_{d1}}{Y_f+K_{d2}}.
\label{eq:r4-2}
\end{equation}
Since $K_{d2}>K_{d1}$, the numerator is positive, so $1-\rho>0$, i.e.,
$\rho<1$.

\textbf{Step 4: $\rho\approx1$.} In \eqref{eq:r4-2}, the numerator
satisfies $K_{d2}-K_{d1}<K_{d2}$, and the denominator satisfies
$Y_f+K_{d2}>Y_f\ge E$ by \eqref{eq:r4-1}. A larger numerator and a
smaller denominator each make a positive fraction larger, so
\begin{equation}
0<1-\rho<\frac{K_{d2}}{Y_f}\le\frac{K_{d2}}{E}=\frac{K_{d2}}{Y-S-(C+D)}\ll1.
\label{eq:r4-3}
\end{equation}
Therefore $\rho\approx1$, and by the same argument for the treatment
condition, $\rho'\approx1$.

\textbf{Step 5: comparing $\rho'$ and $\rho$.} First, $\rho$ increases
with $Y_f$. Differentiating with the quotient rule,
$$
\frac{d\rho}{dY_f}
=\frac{(1)(Y_f+K_{d2})-(Y_f+K_{d1})(1)}{(Y_f+K_{d2})^2}
=\frac{K_{d2}-K_{d1}}{(Y_f+K_{d2})^2}
>0.
$$
Second, $Y_f'>Y_f$ when $S_t<S_c$. To see this, substitute
$\theta_1=Y_f/(Y_f+K_{d1})$ and $\theta_2=Y_f/(Y_f+K_{d2})$ into antibody
conservation, and define the right-hand side as a function of a trial
free-antibody concentration $x\ge 0$ and the target total $S$:
$$
F(x,S)=x+S\frac{x}{x+K_{d1}}+(C+D)\frac{x}{x+K_{d2}}.
$$
The equilibrium $Y_f$ is the value of $x$ that solves $F(x,S)=Y$.
\begin{itemize}
    \item $F$ is increasing in $x$: each term is increasing, since
    $\frac{d}{dx}x=1>0$ and, by the quotient rule, for any $K>0$,
    $$
    \frac{d}{dx}\frac{x}{x+K}
    =\frac{(1)(x+K)-x(1)}{(x+K)^2}
    =\frac{K}{(x+K)^2}>0.
    $$
    \item $F$ is increasing in $S$ for $x>0$: $\partial F/\partial S
    =x/(x+K_{d1})=\theta_1>0$.
\end{itemize}
The two conditions share $Y$, so $F(Y_f,S_c)=Y=F(Y_f',S_t)$. Because
$S_t<S_c$ and $Y_f>0$ (by \eqref{eq:r4-1}), the second property gives
$$
F(Y_f,S_t)<F(Y_f,S_c)=Y=F(Y_f',S_t).
$$
Since $F(\cdot,S_t)$ is increasing, $F(Y_f,S_t)<F(Y_f',S_t)$ requires
$Y_f<Y_f'$. In words: with fewer target epitopes to absorb antibody in
the treatment condition, more antibody is left free.

Combining the two facts, $Y_f'>Y_f$ implies $\rho'>\rho$, just as in
Regime 1. But both lie below 1 (Step 3), so using \eqref{eq:r4-3},
$$
0<\rho'-\rho<1-\rho<\frac{K_{d2}}{E_c}\ll1.
$$
The difference is negligible: $\rho\approx\rho'\approx1$.

\textbf{Result.} The immunoprecipitation contains
$A\theta_1:B\theta_1:C\theta_2:D\theta_2=A:B:\rho C:\rho D$, so
$$
\text{IP composition}\approx A:B:C:D,
$$
which is the same as the input. Because nearly every epitope is captured
(Step 2), the immunoprecipitation is essentially a copy of the
epitope-containing chromatin, with no preference for target over
off-target epitopes.

\subsubsection{ChIP normalization}

We reuse the read model and normalization formulas from Regime 1,
substituting $\rho\approx\rho'\approx1$.

\textbf{Scaling by spike-in reads (ChIP-Rx).} From Regime 1, the ratio
between conditions is
$$
\frac{n_i'/N_s'}{n_i/N_s}
=\frac{a_i'}{a_i}\cdot\frac{B+\rho D}{B+\rho' D}.
$$
Since $\rho'>\rho$ (Step 5), $B+\rho'D>B+\rho D$, so the bias factor is
slightly less than 1. To bound how far below 1 it can be, use
$\rho'<1$ to enlarge the denominator, $B+\rho'D<B+D$, which makes the
fraction smaller:
$$\begin{aligned}
\frac{B+\rho D}{B+\rho' D}
&>\frac{B+\rho D}{B+D}
=\frac{(B+D)-(1-\rho)D}{B+D}
=1-(1-\rho)\frac{D}{B+D}.
\end{aligned}$$
Since $D/(B+D)\le1$ and, by \eqref{eq:r4-3}, $1-\rho<K_{d2}/E_c$,
$$
1-\frac{K_{d2}}{E_c}<\frac{B+\rho D}{B+\rho' D}<1.
$$
So the bias factor is $\approx1$, and \textbf{ChIP-Rx normalization is
essentially unbiased} at target loci in this idealized model: both
conditions capture both spike-in epitope classes with the same (nearly
complete) efficiency, so the spike-in reads track the amount of chromatin
just as they would in the input.

\textbf{Scaling by total reads (RPM).} Substituting $\rho\approx\rho'\approx1$
into the Regime 1 ratio,
$$
\frac{n_i'/N'}{n_i/N}
=\frac{a_i'}{a_i}\cdot\frac{A+B+\rho(C+D)}{A'+B+\rho'(C+D)}
\approx\frac{a_i'}{a_i}\cdot\frac{A+B+C+D}{A'+B+C+D}.
$$
If the off-target pool dominates ($C+D\gg A$), the second factor is
close to 1 and even RPM recovers the global change; the large,
unchanging off-target pool acts as an internal reference. This is not a
virtue, as we now see.

\textbf{The cost: loss of specificity.} Let $c_j$ be the off-target
epitope concentration at an off-target locus $j$. The immunoprecipitation
contains $a_i\theta_1$ at target locus $i$ and $c_j\theta_2$ at locus
$j$. Dividing both by $\theta_1$ and substituting $\rho=\theta_2/\theta_1$,
$$
\frac{a_i\theta_1}{c_j\theta_2}
=\frac{a_i}{\rho\,c_j}
\approx\frac{a_i}{c_j},
$$
which is the same ratio as in the input. Compare Regime 1, where
$\rho\approx r$ boosts target over off-target by a factor $1/r=K_{d2}/K_{d1}\gg1$.
Here that enrichment is lost entirely. Summing over loci, the fraction
of IP reads that come from endogenous off-target epitopes is
$$
\frac{\rho C}{A+B+\rho(C+D)}\approx\frac{C}{A+B+C+D},
$$
and the fraction from all off-target epitopes is approximately
$(C+D)/(A+B+C+D)$. When off-target epitopes are far more abundant than
the target (as is typical, e.g., for an antibody that weakly recognizes
unmodified histones), most reads are off-target and peaks at target loci
are diluted into the background. Normalization is unbiased, but the
signal it normalizes is mostly noise.

\textbf{In practice.} This regime is an idealization. Real ChIP
``background'' includes not only weak antibody-epitope binding (which
saturates, as modeled here) but also chromatin that sticks
nonspecifically to beads, tubes, or antibody-bead complexes. That
background does not follow the binding isotherm $Y_f/(Y_f+K_{d2})$; it
roughly scales with the amount of input chromatin and need not be the
same in the endogenous and spike-in chromatin. Moreover, reaching
$Y>S+C+D+K_{d2}$ requires enormous amounts of antibody when the
off-target pool is large. The main lesson of this regime is therefore
not that saturating antibody makes spike-in normalization safe, but that
adding antibody beyond the point of saturating the target trades
specificity for (at best) uniform capture: the immunoprecipitation
increasingly resembles the input.

\subsection{Crossovers between regimes}

Regimes 1--4 each describe a limit in which one or two terms of antibody
conservation dominate. Real experiments need not sit deep inside any one
limit, and the boundaries between regimes are exactly where the
treatment and control conditions are most likely to behave differently.
In this section we first prove a result that holds in \emph{every}
regime, and then work out the three boundary regions explicitly:
\begin{itemize}
    \item the \textbf{equivalence crossover} ($Y\approx S$), between
    regimes 1 and 3;
    \item the \textbf{off-target half-saturation crossover}
    ($Y_f\sim K_{d2}$), between regimes 3 and 4;
    \item the \textbf{weak-binding crossover} ($K_{d1}\sim S$), between
    regimes 1 and 2.
\end{itemize}
As before, $S\in\{S_c,S_t\}$ with $S_c=A+B$ and $S_t=A'+B$, and
$r=K_{d1}/K_{d2}$. We abbreviate the off-target pool as $M=C+D$, which is
the same in both conditions.

\subsubsection{A general principle: less target means more free antibody}

\textbf{Step 1: $\rho$ depends only on $Y_f$, and increases with it.}
By definition,
$$
\rho=\frac{Y_f+K_{d1}}{Y_f+K_{d2}},
$$
which involves no concentration other than $Y_f$. Treating $\rho$ as a
function of $Y_f$ and differentiating with the quotient rule,
$$
\frac{d\rho}{dY_f}
=\frac{1\cdot(Y_f+K_{d2})-(Y_f+K_{d1})\cdot1}{(Y_f+K_{d2})^2}
=\frac{K_{d2}-K_{d1}}{(Y_f+K_{d2})^2}.
$$
Since $K_{d1}<K_{d2}$, the numerator is positive, and the denominator is a
square, so
\begin{equation}
\frac{d\rho}{dY_f}=\frac{K_{d2}-K_{d1}}{(Y_f+K_{d2})^2}>0.
\label{eq:x-drho}
\end{equation}
Thus, \textbf{whichever condition has more free antibody has the larger
$\rho$}. Intuitively, free antibody is the ``pressure'' that pushes
antibody onto epitopes. Target epitopes, with their small $K_{d1}$, are
already mostly occupied at low pressure, so raising the pressure helps
the weak off-target epitopes more than the target epitopes.

\textbf{Step 2: $Y_f$ is set implicitly by antibody conservation.}
Substituting $\theta_1=Y_f/(Y_f+K_{d1})$ and $\theta_2=Y_f/(Y_f+K_{d2})$
into antibody conservation $Y=Y_f+S\theta_1+M\theta_2$ gives
\begin{equation}
Y=F(Y_f;S),
\qquad
F(y;S)=y+\frac{Sy}{y+K_{d1}}+\frac{My}{y+K_{d2}}.
\label{eq:x-F}
\end{equation}
Here $F(y;S)$ is the total antibody that would be needed to produce a free
antibody concentration $y$ when the total target is $S$: $y$ free, plus
$Sy/(y+K_{d1})$ bound to target, plus $My/(y+K_{d2})$ bound off-target.
We cannot in general solve \eqref{eq:x-F} for $Y_f$ in closed form (after
clearing denominators it is a cubic), but we can learn a lot from how $F$
depends on its two arguments.

\emph{$F$ is increasing in $y$.} Each fraction has the form
$y/(y+K)$, whose derivative by the quotient rule is
$$
\frac{d}{dy}\left(\frac{y}{y+K}\right)
=\frac{1\cdot(y+K)-y\cdot1}{(y+K)^2}
=\frac{K}{(y+K)^2}.
$$
Applying this with $K=K_{d1}$ and $K=K_{d2}$,
\begin{equation}
\frac{\partial F}{\partial y}
=1+\frac{SK_{d1}}{(y+K_{d1})^2}+\frac{MK_{d2}}{(y+K_{d2})^2}
\ge1>0.
\label{eq:x-dFdy}
\end{equation}
(The symbol $\partial$ means we differentiate with respect to $y$ while
holding $S$ fixed.) Also $F(0;S)=0$ and $F(y;S)\ge y$, so $F$ climbs
continuously from $0$ to arbitrarily large values. By the intermediate
value theorem, for any $Y>0$ there is a solution $Y_f$ of
$F(Y_f;S)=Y$, and because $F$ is strictly increasing in $y$, it is the only
one. Moreover, $Y_f>0$, since $F(0;S)=0<Y$.

\emph{$F$ is increasing in $S$.} $S$ appears only in the middle term, so
$$
\frac{\partial F}{\partial S}=\frac{y}{y+K_{d1}}=\theta_1>0
\quad\text{for } y>0.
$$
More antibody is needed to reach a given free concentration when there
is more target to soak it up.

\textbf{Step 3: smaller $S$ gives larger $Y_f$.} Let $Y_f$ and $Y_f'$ be
the solutions for the control ($S_c$) and treatment ($S_t$), with the same
$Y$, and suppose $A'<A$, so $S_t<S_c$. Evaluate the treatment's function
at the control's free antibody. Since only the middle term depends on
$S$,
$$
F(Y_f;S_t)
=F(Y_f;S_c)-\frac{(S_c-S_t)\,Y_f}{Y_f+K_{d1}}
=Y-(S_c-S_t)\,\theta_1,
$$
where we substituted $F(Y_f;S_c)=Y$ from \eqref{eq:x-F}. The subtracted
term is positive (because $S_c>S_t$ and $\theta_1>0$), so
$$
F(Y_f;S_t)<Y=F(Y_f';S_t).
$$
Since $F(\,\cdot\,;S_t)$ is increasing \eqref{eq:x-dFdy}, a smaller output
requires a smaller input, so $Y_f<Y_f'$. In words: removing target
epitopes releases the antibody that would have bound them, and the
system re-equilibrates at a higher free-antibody concentration.

\textbf{Step 4: result.} Combining Step 3 with \eqref{eq:x-drho},
\begin{equation}
A'<A
\quad\Longrightarrow\quad
Y_f'>Y_f
\quad\Longrightarrow\quad
\rho'>\rho.
\label{eq:x-mono}
\end{equation}
This holds in \emph{every} regime, without any approximation. The
regime-specific formulas only tell us \emph{how much} larger $\rho'$ is:
in the deep limits of regimes 2 and 4 the difference is negligible
($\rho'\approx\rho$), whereas near the crossovers below it can be large.

\textbf{How large can the difference be?} The same calculation can be
phrased with implicit differentiation. Differentiating both sides of
$Y=F(Y_f(S);S)$ with respect to $S$ (with $Y$ fixed) using the chain rule,
$$
0=\frac{\partial F}{\partial y}\frac{dY_f}{dS}+\frac{\partial F}{\partial S},
\qquad\text{so}\qquad
\frac{dY_f}{dS}
=-\frac{\partial F/\partial S}{\partial F/\partial y}
=-\frac{\theta_1}{\partial F/\partial y}.
$$
Since $0<\theta_1\le1$ and $\partial F/\partial y\ge1$ by
\eqref{eq:x-dFdy}, we get $-1\le dY_f/dS<0$. By the mean value theorem,
there is some $S^*$ between $S_t$ and $S_c$ with
$$
Y_f'-Y_f
=-\frac{dY_f}{dS}\bigg|_{S^*}\,(S_c-S_t)
\le S_c-S_t
=A-A'.
$$
That is, removing $A-A'$ target epitopes can raise the free antibody by at
most $A-A'$ (the released antibody can only be shared out among solution,
the remaining target, and off-target sites). Applying the mean value
theorem again, now to $\rho$ as a function of $Y_f$, and bounding the
derivative \eqref{eq:x-drho} by its largest value (at $Y_f=0$),
\begin{equation}
\rho'-\rho
\le\frac{K_{d2}-K_{d1}}{K_{d2}^2}\,(Y_f'-Y_f)
\le\frac{(1-r)(A-A')}{K_{d2}}
<\frac{A-A'}{K_{d2}}.
\label{eq:x-bound}
\end{equation}
So the \emph{absolute} change in $\rho$ is always modest unless
$A-A'$ is comparable to $K_{d2}$. The \emph{relative} change
$\rho'/\rho$ can nevertheless be large when $\rho$ itself is small
($\rho\sim r$) and $Y_f$ is very sensitive to $S$, which is exactly the
situation near equivalence.

\subsubsection{Equivalence crossover (between regimes 1 and 3)}

Regime 1 required $\sqrt{K_{d1}Y}\ll S-Y$, and regime 3 requires the
antibody to exceed the target. Between them lies a window around
$Y\approx S$ where antibody and target \emph{titrate} each other. Define
the excess target
$$
\delta=S-Y,
$$
which is positive on the regime-1 side and negative on the regime-3 side.

\textbf{Assumptions.}
\begin{enumerate}
    \item Off-target-bound antibody is negligible compared with free
    antibody: $M\theta_2\ll Y_f$. Since
    $M\theta_2=MY_f/(Y_f+K_{d2})\le MY_f/K_{d2}$, a sufficient condition
    is $M\ll K_{d2}$. (We relax this in a remark below.)
    \item Target binding is tight compared with the amount of target:
    $K_{d1}\ll S$.
    \item We are near equivalence: $|\delta|\ll S$.
    \item Off-target epitopes are far from half-occupied: $Y_f\ll K_{d2}$.
\end{enumerate}

\textbf{Step 1: an exact quadratic.} Dropping $M\theta_2$ (assumption 1)
from \eqref{eq:r1-1}, the free target is
$$
S(1-\theta_1)=(S-Y)+Y_f=\delta+Y_f,
$$
and solving antibody conservation for the bound target gives
$S\theta_1=Y-Y_f$. The mass-action relation for target,
$K_{d1}=Y_f(A_f+B_f)/(Y_A+Y_B)$, has free target $A_f+B_f=S(1-\theta_1)$
and bound target $Y_A+Y_B=S\theta_1$. Substituting both and multiplying
through by $Y-Y_f$,
$$
K_{d1}(Y-Y_f)=Y_f(\delta+Y_f).
$$
Expanding and collecting everything on one side,
\begin{equation}
Y_f^2+(\delta+K_{d1})Y_f-K_{d1}Y=0.
\label{eq:x-quad-exact}
\end{equation}
Given assumption 1, this is exact; we will reuse it for the weak-binding
crossover.

\textbf{Step 2: simplify near equivalence.} Two terms in
\eqref{eq:x-quad-exact} are small corrections near equivalence. First,
substituting $Y=S-\delta$,
$$
K_{d1}Y=K_{d1}S\left(1-\frac{\delta}{S}\right)\approx K_{d1}S
$$
by assumption 3. Second, comparing $K_{d1}Y_f$ with $K_{d1}Y$, their ratio
is $Y_f/Y$, which is $\ll1$ provided $Y_f\ll Y\approx S$ (we verify this in
Step 4). Dropping $K_{d1}Y_f$ and replacing $K_{d1}Y$ by $K_{d1}S$,
\begin{equation}
Y_f\,(Y_f+\delta)\approx K_{d1}S.
\label{eq:x-quad}
\end{equation}
This has a direct interpretation. $Y_f+\delta$ is the free target, and
nearly all of the $\approx S$ target epitopes are bound, so
\eqref{eq:x-quad} is just mass action:
(free antibody)$\times$(free target) $=K_{d1}\times$(bound target).

\textbf{Step 3: solve the quadratic.} Writing \eqref{eq:x-quad} as
$Y_f^2+\delta Y_f-K_{d1}S=0$, the quadratic formula gives
$$
Y_f=\frac{-\delta\pm\sqrt{\delta^2+4K_{d1}S}}{2}.
$$
The product of the two roots equals the constant term, $-K_{d1}S<0$, so one
root is positive and one is negative. Free antibody must be positive, so
we take the $+$ sign. Multiplying the numerator and denominator by the
conjugate $\delta+\sqrt{\delta^2+4K_{d1}S}$ and using
$(\sqrt{\delta^2+4K_{d1}S})^2-\delta^2=4K_{d1}S$ gives an equivalent form
that is convenient when $\delta>0$:
\begin{equation}
Y_f
=\frac{-\delta+\sqrt{\delta^2+4K_{d1}S}}{2}
=\frac{2K_{d1}S}{\delta+\sqrt{\delta^2+4K_{d1}S}}.
\label{eq:x-yf}
\end{equation}
The natural scale is $\sqrt{K_{d1}S}$. Writing
$\delta=x\sqrt{K_{d1}S}$ and dividing \eqref{eq:x-yf} by $\sqrt{K_{d1}S}$,
\begin{equation}
\frac{Y_f}{\sqrt{K_{d1}S}}
=h(x)
=\frac{-x+\sqrt{x^2+4}}{2}
=\frac{2}{x+\sqrt{x^2+4}}.
\label{eq:x-h}
\end{equation}
Comparing the two forms of $h$, $h(-x)=(x+\sqrt{x^2+4})/2=1/h(x)$: excess
target and excess antibody of the same size shift $Y_f$ by the same factor
in opposite directions.

\textbf{Step 4: the three limits.}
\begin{itemize}
    \item \emph{Excess target, $\delta\gg\sqrt{K_{d1}S}$ (regime 1).} Then
    $4K_{d1}S\ll\delta^2$, so $\sqrt{\delta^2+4K_{d1}S}\approx\delta$.
    Substituting into the second form of \eqref{eq:x-yf},
    $$
    Y_f\approx\frac{2K_{d1}S}{\delta+\delta}=\frac{K_{d1}S}{\delta}
    =\frac{K_{d1}S}{S-Y}.
    $$
    Since $S\approx Y$ (assumption 3), this matches the regime-1 estimate
    $K_{d1}Y/(S-Y)$ of \eqref{eq:r1-4}. (The exact quadratic
    \eqref{eq:x-quad-exact} gives $K_{d1}Y/(S-Y+K_{d1})$ in this limit,
    matching even more closely.)
    \item \emph{Exact equivalence, $\delta=0$.} Substituting
    $\delta=0$ into \eqref{eq:x-yf},
    $$
    Y_f=\frac{\sqrt{4K_{d1}S}}{2}=\sqrt{K_{d1}S},
    $$
    the geometric mean of $K_{d1}$ and $S$, as anticipated in the
    discussion of regime 1's assumption 2.
    \item \emph{Excess antibody, $-\delta\gg\sqrt{K_{d1}S}$ (regime 3).}
    Now $\delta=-|\delta|$ and again $\sqrt{\delta^2+4K_{d1}S}\approx|\delta|$.
    Substituting into the first form of \eqref{eq:x-yf},
    $$
    Y_f\approx\frac{|\delta|+|\delta|}{2}=|\delta|=Y-S.
    $$
    The target is saturated, and the excess antibody is free (it has
    nowhere else to go, since off-target binding is negligible).
\end{itemize}
We can now verify the claim $Y_f\ll S$ used in Step 2. Using
$\sqrt{a+b}\le\sqrt a+\sqrt b$ for $a,b\ge0$ in the first form of
\eqref{eq:x-yf},
$$
Y_f\le\frac{-\delta+|\delta|+2\sqrt{K_{d1}S}}{2}\le|\delta|+\sqrt{K_{d1}S}.
$$
Here $|\delta|\ll S$ by assumption 3, and $\sqrt{K_{d1}S}\ll S$ because,
dividing by $\sqrt S$, it is equivalent to $\sqrt{K_{d1}}\ll\sqrt S$, i.e.,
assumption 2. So $Y_f\ll S$.

\textbf{Step 5: $\rho$.} By assumption 4, $Y_f+K_{d2}\approx K_{d2}$, so
\begin{equation}
\rho
=\frac{Y_f+K_{d1}}{Y_f+K_{d2}}
\approx\frac{Y_f+K_{d1}}{K_{d2}}
=r\left(1+\frac{Y_f}{K_{d1}}\right),
\label{eq:x-rho-eq}
\end{equation}
where the last step factors out $K_{d1}/K_{d2}=r$. Substituting the three
limits of Step 4:
$$
\rho\approx
\begin{cases}
r\left(1+\dfrac{S}{S-Y}\right)\approx r\dfrac{S}{S-Y}, & S-Y\gg\sqrt{K_{d1}S}
\quad\text{(regime 1)},\\[2ex]
r\left(1+\sqrt{\dfrac{S}{K_{d1}}}\right)\approx\dfrac{\sqrt{K_{d1}S}}{K_{d2}}, & Y=S,\\[2ex]
r+\dfrac{Y-S}{K_{d2}}, & Y-S\gg\sqrt{K_{d1}S}
\quad\text{(regime 3, }M\ll K_{d2}\text{)}.
\end{cases}
$$
(In the first line, $S/(S-Y)\gg1$ near equivalence, so the $1$ is
negligible; in the second, $\sqrt{S/K_{d1}}\gg1$ by assumption 2.)
Through the window, $\rho$ climbs from about $r\,S/(S-Y)$ to about
$(Y-S)/K_{d2}$, and the climb happens over a range of $Y$ of width a few
times $\sqrt{K_{d1}S}$, i.e., a \emph{relative} width of order
$\sqrt{K_{d1}S}/S=\sqrt{K_{d1}/S}\ll1$. For example, by \eqref{eq:x-h},
moving from $x=2$ to $x=-2$ multiplies $Y_f$ by
$h(-2)/h(2)=1/h(2)^2=(1+\sqrt2)^2\approx5.8$; moving from $x=4$ to $x=-4$
multiplies it by $\bigl((4+\sqrt{20})/2\bigr)^2\approx18$.

\textbf{$\rho'$ versus $\rho$.} The treatment has
$\delta_t=S_t-Y=\delta_c-(A-A')<\delta_c$. In \eqref{eq:x-yf} (second
form), the denominator $\delta+\sqrt{\delta^2+4K_{d1}S}$ is increasing in
$\delta$, since its derivative is
$1+\delta/\sqrt{\delta^2+4K_{d1}S}>0$ (because
$|\delta|<\sqrt{\delta^2+4K_{d1}S}$). So decreasing $\delta$ increases
$Y_f$, and by \eqref{eq:x-rho-eq}, $\rho'>\rho$, as guaranteed in general by
\eqref{eq:x-mono}. Because $\rho\propto Y_f+K_{d1}$ here, the ratio is
$$
\frac{\rho'}{\rho}\approx\frac{Y_f'+K_{d1}}{Y_f+K_{d1}},
$$
which is a ratio of free-antibody concentrations. If the drop $A-A'$ is
large enough to carry the treatment across the window while the control
stays on the regime-1 side, this ratio is approximately
$(Y-S_t)\big/\bigl(K_{d1}S_c/(S_c-Y)\bigr)$, which can be enormous because
$K_{d1}$ is small. Equivalently, bound \eqref{eq:x-bound} gives
$\rho'/\rho\le1+(A-A')/(Y_f+K_{d1})$, which is large when $A-A'$ exceeds
the free antibody $\sim\sqrt{K_{d1}S}$.

\textbf{Numerical illustration.} Take $K_{d1}=0.1$~nM, $K_{d2}=1000$~nM
($r=10^{-4}$), $A=90$~nM, $B=10$~nM, $M\ll K_{d2}$, and a treatment that
removes about 11\% of the target, $A'=80$~nM. Then $S_c=100$~nM and
$S_t=90$~nM, and $\sqrt{K_{d1}S}\approx3$~nM.
\begin{itemize}
    \item With $Y=95$~nM, $\delta_c=5$~nM and $\delta_t=-5$~nM. Substituting
    into \eqref{eq:x-yf},
    $$
    \begin{gathered}
    Y_f\approx\frac{-5+\sqrt{25+4(0.1)(100)}}{2}=\frac{-5+8.06}{2}\approx1.53~\text{nM},\\
    Y_f'\approx\frac{5+\sqrt{25+4(0.1)(90)}}{2}=\frac{5+7.81}{2}\approx6.41~\text{nM}.
    \end{gathered}
    $$
    By \eqref{eq:x-rho-eq}, $\rho\approx1.63/1000=1.6\times10^{-3}$ and
    $\rho'\approx6.51/1000=6.5\times10^{-3}$, so $\rho'/\rho\approx4$.
    (Solving \eqref{eq:x-F} numerically gives $\rho'/\rho\approx4.2$.)
    \item With $Y=50$~nM (deep in regime 1), the same 11\% drop gives
    $\rho'/\rho\approx\dfrac{90/40}{100/50}=1.125$.
\end{itemize}
An 11\% change in target abundance produces a fourfold change in the
relative off-target capture efficiency when the antibody is titrated near
equivalence, but only a 12.5\% change when the antibody is at half the
target concentration.

\textbf{Remark: an off-target sink.} If $M$ is not small compared with
$K_{d2}$ but off-target sites are still unsaturated ($Y_f\ll K_{d2}$), then
$M\theta_2\approx MY_f/K_{d2}$, and the antibody not on target is
$$
U=Y_f+M\theta_2\approx gY_f,
\qquad
g=1+\frac{M}{K_{d2}}.
$$
Repeating Step 1 with $U$ in place of $Y_f$ in \eqref{eq:r1-1} and
in antibody conservation, the free target is $\delta+U$ and the bound target
is $Y-U$, so mass action reads $K_{d1}(Y-U)=Y_f(\delta+U)$. Multiplying both
sides by $g$ and substituting $gY_f=U$,
$$
gK_{d1}(Y-U)=U(\delta+U).
$$
This is \eqref{eq:x-quad-exact} with $Y_f\to U$ and $K_{d1}\to gK_{d1}$.
All the results above carry over with these replacements; in particular,
$U\approx\sqrt{gK_{d1}S}$ at equivalence and $Y_f=U/g$. The off-target pool
acts as a buffer that soaks up a fixed fraction $1-1/g$ of the antibody not
on target, and it widens the window to $\sqrt{gK_{d1}S}$, which makes the
crossover less sharp.

\subsubsection{Off-target half-saturation crossover (between regimes 3 and 4)}

Once $Y-S\gg\sqrt{K_{d1}S}$, the target is saturated and the remaining
question is how the excess antibody distributes between solution and
off-target epitopes. Define the excess antibody
$$
E=Y-S>0.
$$

\textbf{Assumption.} The target is saturated: the free target
$S(1-\theta_1)=SK_{d1}/(Y_f+K_{d1})$ is negligible compared with $E$. In this
crossover $Y_f\sim K_{d2}\gg K_{d1}$, so the free target is about
$SK_{d1}/K_{d2}=rS$, and the assumption amounts to $rS\ll Y-S$.

\textbf{Step 1: a quadratic for $Y_f$.} Saturation means
$S\theta_1=S-S(1-\theta_1)\approx S$. Substituting this and
$\theta_2=Y_f/(Y_f+K_{d2})$ into antibody conservation,
$Y=Y_f+S+MY_f/(Y_f+K_{d2})$, and subtracting $S$ from both sides,
$$
E=Y_f+\frac{MY_f}{Y_f+K_{d2}}.
$$
In words, the excess antibody is either free or bound off-target.
Multiplying both sides by $Y_f+K_{d2}$,
$$
EY_f+EK_{d2}=Y_f^2+K_{d2}Y_f+MY_f,
$$
and collecting terms on one side,
\begin{equation}
Y_f^2+bY_f-EK_{d2}=0,
\qquad
b=K_{d2}+M-E.
\label{eq:x-quad-c}
\end{equation}

\textbf{Step 2: solve.} The product of the roots is $-EK_{d2}<0$, so there
is exactly one positive root. By the quadratic formula, and rationalizing
as in \eqref{eq:x-yf},
\begin{equation}
Y_f
=\frac{-b+\sqrt{b^2+4EK_{d2}}}{2}
=\frac{2EK_{d2}}{b+\sqrt{b^2+4EK_{d2}}},
\qquad
\rho=\frac{Y_f+K_{d1}}{Y_f+K_{d2}}\approx\frac{Y_f}{Y_f+K_{d2}},
\label{eq:x-yf-c}
\end{equation}
where the last approximation uses $K_{d1}\ll Y_f$.

\textbf{Step 3: the limits.}
\begin{itemize}
    \item \emph{Small excess, $E\ll K_{d2}+M$ (regime 3).} Then
    $b\approx K_{d2}+M>0$ and $4EK_{d2}\le4E(K_{d2}+M)\ll b^2$, so
    $\sqrt{b^2+4EK_{d2}}\approx b$. Substituting into the second form of
    \eqref{eq:x-yf-c},
    $$
    Y_f\approx\frac{2EK_{d2}}{2b}\approx\frac{K_{d2}(Y-S)}{K_{d2}+M}.
    $$
    This is $\ll K_{d2}$, so $\rho\approx(Y_f+K_{d1})/K_{d2}$; substituting
    $Y_f$ gives $\rho\approx r+(Y-S)/(K_{d2}+C+D)$, the regime-3 result.
    \item \emph{Large excess, $E\gg K_{d2}+M$ (regime 4).} Now
    $b=-(E-M-K_{d2})$ is large and negative. Factoring out $|b|$ and using
    the linear approximation $\sqrt{1+u}\approx1+u/2$ for small $u$,
    $$
    \sqrt{b^2+4EK_{d2}}
    =|b|\sqrt{1+\frac{4EK_{d2}}{b^2}}
    \approx|b|+\frac{2EK_{d2}}{|b|}
    \approx|b|+2K_{d2},
    $$
    where the last step uses $E/|b|\approx1$. Substituting into the first form
    of \eqref{eq:x-yf-c},
    $$
    Y_f\approx\frac{|b|+|b|+2K_{d2}}{2}=|b|+K_{d2}=E-M=Y-S-(C+D).
    $$
    All epitopes are saturated and the leftover antibody is free. Then
    $\rho\approx(E-M)/(E-M+K_{d2})\to1$ when $E-M\gg K_{d2}$, the regime-4
    result.
    \item \emph{Half-saturation.} $\rho\approx1/2$ when $Y_f=K_{d2}$, i.e.,
    when off-target epitopes are half occupied. Rather than solving
    \eqref{eq:x-quad-c}, substitute $Y_f=K_{d2}$ directly into
    $E=Y_f+MY_f/(Y_f+K_{d2})$:
    $$
    E_{1/2}=K_{d2}+\frac{M}{2},
    \qquad\text{i.e.,}\qquad
    Y=S+K_{d2}+\frac{C+D}{2}.
    $$
    The antibody must cover the target ($S$), half of the off-target pool
    ($M/2$), and leave $K_{d2}$ free.
\end{itemize}

\textbf{Two shapes of crossover.} The size of $M$ relative to $K_{d2}$
controls how the transition looks.
\begin{itemize}
    \item If $M\ll K_{d2}$, set $M=0$ in \eqref{eq:x-quad-c}:
    $Y_f^2+(K_{d2}-E)Y_f-EK_{d2}=0$. This factors as
    $(Y_f-E)(Y_f+K_{d2})=0$ (check by expanding:
    $Y_f^2+K_{d2}Y_f-EY_f-EK_{d2}$), so $Y_f=E$ and
    $$
    \rho\approx\frac{(Y-S)+K_{d1}}{(Y-S)+K_{d2}}.
    $$
    This is a gentle hyperbola: $\rho$ rises from $0.09$ to $0.91$ as
    $Y-S$ goes from $0.1K_{d2}$ to $10K_{d2}$, spanning two orders of
    magnitude.
    \item If $M\gg K_{d2}$, then \eqref{eq:x-quad-c} has exactly the form of
    \eqref{eq:x-quad-exact} with $Y_f$ unchanged, $K_{d1}\to K_{d2}$,
    $\delta\to M-E$, and $Y\to E$. So the antibody titrates the off-target
    pool the way it titrated the target in the equivalence crossover: there
    is a ``second equivalence'' at $E\approx M$ with window width
    $\sqrt{K_{d2}M}$. Half-saturation ($E=K_{d2}+M/2$) occurs before this
    second equivalence, and by $E=M$ we already have
    $Y_f\approx\sqrt{K_{d2}M}\gg K_{d2}$, so $\rho$ is close to 1.
\end{itemize}

\textbf{$\rho'$ versus $\rho$.} Because the target is saturated, $S$ enters
only through $E=Y-S$, and the treatment has
$E'=E+(A-A')>E$. Differentiating $E=Y_f+MY_f/(Y_f+K_{d2})$ implicitly with
respect to $E$,
$$
1=\frac{dY_f}{dE}\left(1+\frac{MK_{d2}}{(Y_f+K_{d2})^2}\right),
\qquad\text{so}\qquad
0<\frac{dY_f}{dE}\le1.
$$
Thus $Y_f'>Y_f$ and $\rho'>\rho$, as \eqref{eq:x-mono} requires. But the
effect is small. Near half-saturation, \eqref{eq:x-drho} with $Y_f=K_{d2}$
gives $d\rho/dY_f\approx K_{d2}/(2K_{d2})^2=1/(4K_{d2})$, so by the chain rule
$$
\rho'-\rho
\approx\frac{d\rho}{dY_f}\cdot\frac{dY_f}{dE}\cdot(A-A')
\le\frac{A-A'}{4K_{d2}},
\qquad
\frac{\rho'}{\rho}\lesssim1+\frac{A-A'}{2K_{d2}},
$$
using $\rho\approx1/2$ in the second expression. Unlike near equivalence,
$\rho$ is not small here, so a small absolute change in $\rho$ cannot become
a large relative change. With the parameters of the earlier illustration
and $C=D=1000$~nM, taking $Y=2500$~nM gives (numerically) $\rho\approx0.561$
and $\rho'\approx0.562$.

\subsubsection{Weak-binding crossover (between regimes 1 and 2)}

When $K_{d1}$ is comparable to $S$, neither ``all target bound'' nor
``almost no antibody bound'' holds. Clearing the denominators of
\eqref{eq:x-F} by multiplying both sides by $(Y_f+K_{d1})(Y_f+K_{d2})$,
$$
Y(Y_f+K_{d1})(Y_f+K_{d2})
=Y_f(Y_f+K_{d1})(Y_f+K_{d2})+SY_f(Y_f+K_{d2})+MY_f(Y_f+K_{d1}),
$$
which is a cubic in $Y_f$. Its formula is not illuminating, so in practice
one solves it numerically (Step 2 of the general principle guarantees a
unique positive solution). Two observations remove most of the need to do
so.

First, if off-target-bound antibody is negligible ($M\ll K_{d2}$), the
derivation of \eqref{eq:x-quad-exact} never used $K_{d1}\ll S$, so with
$\delta=S-Y$,
$$
Y_f=\frac{-(S-Y+K_{d1})+\sqrt{(S-Y+K_{d1})^2+4K_{d1}Y}}{2},
\qquad
\rho=\frac{Y_f+K_{d1}}{Y_f+K_{d2}}
$$
is exact across the whole crossover. (If $M$ is not negligible but
off-target sites are unsaturated, the remark on the off-target sink applies.)

Second, the regime-2 limit follows directly from conservation. Since
$\theta_1=Y_f/(Y_f+K_{d1})\le Y_f/K_{d1}$ and $Y_f\le Y$, the antibody bound
to target satisfies
$$
Y-Y_f\approx S\theta_1\le\frac{SY_f}{K_{d1}}\le\frac{S}{K_{d1}}\,Y,
$$
which is $\ll Y$ when $S\ll K_{d1}$, giving $Y_f\approx Y$. As $S/K_{d1}$
grows toward 1, a growing fraction of antibody is captured by the target,
$Y_f$ drops below $Y$, and $\rho$ drops from $(Y+K_{d1})/(Y+K_{d2})$ toward
regime 1's $r\,S/(S-Y)$ (if $Y<S$) or toward the equivalence or regime-3
behavior (if $Y\gtrsim S$). Throughout, \eqref{eq:x-mono} still guarantees
$\rho'>\rho$ when $A'<A$. The difference is small near regime 2: from the
implicit derivative in the general principle, $|dY_f/dS|\le\theta_1\le
Y_f/K_{d1}$, which is small when $Y_f\ll K_{d1}$.

\subsubsection{ChIP normalization}

From Regime 1's ChIP normalization, scaling by total spike-in reads
(ChIP-Rx) gives
$$
\frac{n_i'/N_s'}{n_i/N_s}
=\frac{a_i'}{a_i}\cdot\beta,
\qquad
\beta=\frac{B+\rho D}{B+\rho' D}.
$$
That derivation used only the read model, not any regime-specific
approximation for $\rho$, so it holds everywhere.

\textbf{The bias never favors the treatment.} By \eqref{eq:x-mono},
$\rho'>\rho$ whenever $A'<A$. Multiplying by $D\ge0$ and adding $B$,
$$
B+\rho' D\ge B+\rho D>0,
$$
so dividing both sides by $B+\rho'D$ gives $\beta\le1$, with equality only
if $D=0$ (or approximately, in the deep limits where $\rho'\approx\rho$). The
measured fold change is the true fold change multiplied by $\beta\le1$.
For a global decrease, $a_i'/a_i<1$, so the measured value is even smaller:
\textbf{ChIP-Rx with total spike-in reads never understates, and generally
overstates, a global decrease}, in every regime. (Equivalently, every locus
in the treatment is biased downward by the same factor $\beta$.)

\textbf{How large can the bias be?} Compare $\beta$ with $\rho/\rho'$.
Putting both over a common denominator,
$$
\beta-\frac{\rho}{\rho'}
=\frac{\rho'(B+\rho D)-\rho(B+\rho' D)}{\rho'(B+\rho' D)}
=\frac{\rho' B+\rho\rho' D-\rho B-\rho\rho' D}{\rho'(B+\rho' D)}
=\frac{B(\rho'-\rho)}{\rho'(B+\rho' D)}
\ge0.
$$
Therefore
$$
\frac{\rho}{\rho'}\le\beta\le1.
$$
The bias approaches its worst case $\rho/\rho'$ when $\rho'D\gg B$, i.e.,
when most spike-in reads come from off-target spike-in epitopes, and
vanishes when $\rho'D\ll B$. Combined with the results above:
\begin{itemize}
    \item Deep in regimes 1 (with $Y\ll S$), 2, and 4, $\rho'\approx\rho$,
    so $\beta\approx1$.
    \item Near equivalence, $\rho'/\rho$ can be large (fourfold in the
    illustration), so $\beta$ can be far below 1 if $\rho'D$ is comparable
    to or larger than $B$. This is where the bias is largest.
    \item Near off-target half-saturation, $\rho'/\rho\lesssim1+(A-A')/(2K_{d2})$,
    so $\beta$ is close to 1 unless $A-A'$ is comparable to $K_{d2}$.
\end{itemize}

\textbf{Numerical illustration.} Use $K_{d1}=0.1$~nM, $K_{d2}=1000$~nM,
$A=90$~nM, $A'=80$~nM, $B=10$~nM, and now a substantial off-target pool
$C=D=1000$~nM (so $g=3$ in the off-target-sink remark). The true global fold
change is $80/90\approx0.89$. Solving \eqref{eq:x-F} numerically:
\begin{itemize}
    \item $Y=50$~nM (regime 1): $\beta\approx0.998$.
    \item $Y=95$~nM (near equivalence): $\rho\approx1.20\times10^{-3}$,
    $\rho'\approx2.82\times10^{-3}$, so $\rho D\approx1.2$~nM and
    $\rho'D\approx2.8$~nM, giving $\beta\approx11.2/12.8\approx0.87$. The
    measured fold change is $0.89\times0.87\approx0.78$. Scanning $Y$, the
    smallest $\beta$ ($\approx0.83$) occurs at $Y\approx107$~nM, just past
    equivalence (the off-target sink shifts it slightly).
    \item $Y=2500$~nM (off-target half-saturation): $\beta\approx0.998$.
\end{itemize}
As in Regime 1, counting only reads at the spike-in's target loci
($N_{s,B}$) removes the dependence on $\rho$ and $\rho'$ entirely, and
therefore removes this bias in every regime.






% $$
% A + B = A_f + B_f + Y_A + Y_B
% $$

% Treat $A_f + B_f$ and $Y_A + Y_B$ as unknowns, but $Y_f$, $K_{d1}$, and $A + B$ as knowns. Then,

% $$
% \begin{aligned}
% K_{d1}
% &= \frac{(Y_f) (A_f + B_f)}{Y_A + Y_B} \\
% &= \frac{(Y_f) (A + B - Y_A + Y_B)}{Y_A + Y_B}
% \end{aligned}$$

\end{document}
